\chapter{Probability and Information}\label{ch:info}

Quantum mechanics is a theory of \emph{probabilities}. The only thing Postulate~3 ever promises is that outcome $m$ occurs with probability $p(m)$; it never promises which outcome we will see. To use the theory seriously we must therefore be fluent in probability, and in the two
ideas that turn probabilities into quantitative statements about knowledge and disorder: Fourier analysis, which connects a wave function to its momentum content and underlies the quantum Fourier transform, and the notion of \emph{entropy}, which measures how much we do not know about a state.

This chapter is a self-contained review of those tools, followed by their quantum versions. The route is the following.

\begin{steps}
  \item \textbf{Probability measures} (Section~\ref{sec:prob}) give a precise meaning to ``probability'', both for finitely many outcomes (a qubit measurement) and for a continuum (the position of a particle). We see that every quantum measurement defines one.
  \item \textbf{Fourier transforms} (Section~\ref{sec:fourier}) pass between position and momentum, and between a probability distribution and its characteristic function. Their discrete version is the matrix at the heart of Shor's algorithm.
  \item \textbf{Shannon entropy} (Section~\ref{sec:shannon}) quantifies the information in a classical random variable, and what is needed to store and send it.
  \item \textbf{Von Neumann entropy} (Section~\ref{sec:vonneumann}) extends it to density operators, and measures entanglement.
  \item \textbf{Relative entropy} (Section~\ref{sec:relent}) measures how distinguishable two states are, and gives the cleanest proofs of the basic inequalities.
\end{steps}

\begin{objectives}[In this chapter]
  You will use probability measures, random variables, conditioning and independence; compute Fourier series and transforms and the discrete Fourier transform; compute Shannon and von Neumann entropies and relative entropies; and prove
  their basic properties, including Gibbs' and Klein's inequalities and the subadditivity of the von Neumann entropy.
\end{objectives}

\begin{whereitsits}
  Everything here is built on Chapters~\ref{ch:tools} and~\ref{ch:postulates}: the spectral theorem and functions of operators define the von Neumann entropy, density operators and partial traces
  supply its arguments, and the Schmidt decomposition explains why the entropy of entanglement is well defined. The outcome probabilities of Postulate~3 are the bridge between the two worlds.
\end{whereitsits}

%======================================================================
\section{Probability measures}\label{sec:prob}
%======================================================================

\begin{objectives}
  You will state the axioms of a probability measure and derive its basic properties; work with discrete and continuous distributions, random variables, expectation and variance;
  use conditioning, independence and Bayes' rule; and see why repeated experiments estimate probabilities.
\end{objectives}

\begin{recall}
  Postulate~3 and the outcome probabilities $p(m)$ (Section~\ref{sec:measurement}); orthogonal projectors (Section~\ref{sec:operators}); Bloch vectors (Section~\ref{sec:state-evolution}); Bell states (Section~\ref{sec:composite}).
\end{recall}

Until now probabilities came from a finite list: a qubit measurement has two outcomes, and $p(0)+p(1)=1$. For a position measurement there are uncountably many outcomes and the probability of any single point is zero, so
``a list of numbers'' is no longer enough. The right framework is that of \emph{measure theory}, which handles both cases at once. We need only its basic vocabulary.

\subsection{Probability spaces}

A probability model has three ingredients: the set $\Omega$ of all possible outcomes, a family of subsets of $\Omega$ to which we are willing to assign a probability (the \emph{events}), and the assignment itself.

\begin{definition}[$\sigma$-algebra]\label{def:sigma}
  A \term{$\sigma$-algebra} on a set $\Omega$ is a family $\mathcal F$ of subsets of $\Omega$ such that
  \begin{steps}
    \item $\Omega\in\mathcal F$;
    \item if $A\in\mathcal F$ then its complement $A^c=\Omega\setminus A$ is in $\mathcal F$;
    \item if $A_1,A_2,\dots\in\mathcal F$ then $\bigcup_iA_i\in\mathcal F$ (closed under countable unions).
  \end{steps}
  The elements of $\mathcal F$ are called \term{events}.
\end{definition}

For a finite or countable $\Omega$ we always take $\mathcal F$ to be \emph{all} subsets. For $\Omega=\R$ we take the \term{Borel $\sigma$-algebra}, the smallest one containing all intervals; it contains every set that one meets in practice, and $\sigma$-algebras exist precisely to exclude pathological sets.

\begin{definition}[Probability measure]\label{def:probmeasure}
  A \term{probability measure} on $(\Omega,\mathcal F)$ is a function $\PP:\mathcal F\to[0,1]$ such that
  \begin{steps}
    \item $\PP(\Omega)=1$;
    \item \emph{(countable additivity)} if $A_1,A_2,\dots\in\mathcal F$ are pairwise disjoint, then $\PP\bigl(\bigcup_iA_i\bigr)=\sum_i\PP(A_i)$.
  \end{steps}
  The triple $(\Omega,\mathcal F,\PP)$ is a \term{probability space}.
\end{definition}

These two axioms, due to Kolmogorov, are all there is. Everything else is a consequence.

\begin{proposition}[Basic properties]\label{prop:probprops}
  Let $\PP$ be a probability measure and $A,B,A_1,A_2,\dots$ events. Then
  \begin{enumerate}
    \item[(a)] $\PP(\varnothing)=0$;
    \item[(b)] \emph{finite additivity:} if $A\cap B=\varnothing$ then $\PP(A\cup B)=\PP(A)+\PP(B)$;
    \item[(c)] $\PP(A^c)=1-\PP(A)$;
    \item[(d)] \emph{monotonicity:} if $A\subseteq B$ then $\PP(A)\le\PP(B)$;
    \item[(e)] $\PP(A\cup B)=\PP(A)+\PP(B)-\PP(A\cap B)$;
    \item[(f)] \emph{union bound:} $\PP\bigl(\bigcup_iA_i\bigr)\le\sum_i\PP(A_i)$;
    \item[(g)] \emph{continuity:} if $A_1\subseteq A_2\subseteq\cdots$ and $A=\bigcup_nA_n$, then $\PP(A_n)\to\PP(A)$.
  \end{enumerate}
\end{proposition}
\begin{proof}
  \begin{steps}
    \item (a) Apply countable additivity to $\varnothing,\varnothing,\dots$: $\PP(\varnothing)=\sum_{i=1}^\infty\PP(\varnothing)$. A number in $[0,1]$ equal to an infinite sum of copies of itself must be $0$.
    \item (b) Apply additivity to $A,B,\varnothing,\varnothing,\dots$ and use (a).
    \item (c) $A$ and $A^c$ are disjoint with union $\Omega$, so by (b) $1=\PP(A)+\PP(A^c)$.
    \item (d) Write $B=A\cup(B\setminus A)$, a disjoint union. By (b), $\PP(B)=\PP(A)+\PP(B\setminus A)\ge\PP(A)$.
    \item (e) Decompose into disjoint pieces: $A\cup B=A\cup(B\setminus A)$ and $B=(B\setminus A)\cup(A\cap B)$. Then $\PP(A\cup B)=\PP(A)+\PP(B\setminus A)$ and $\PP(B)=\PP(B\setminus A)+\PP(A\cap B)$; subtract.
    \item (f) For two events, (e) gives $\PP(A\cup B)\le\PP(A)+\PP(B)$; induction gives the finite case. For countably many, define disjoint events $D_1=A_1$, $D_n=A_n\setminus(A_1\cup\dots\cup A_{n-1})$ with the same union and $D_n\subseteq A_n$, and use additivity and (d).
    \item (g) The sets $D_1=A_1$, $D_n=A_n\setminus A_{n-1}$ are disjoint with union $A$, and $A_n=D_1\cup\dots\cup D_n$. So $\PP(A)=\sum_{i}\PP(D_i)=\lim_n\sum_{i\le n}\PP(D_i)=\lim_n\PP(A_n)$. \qedhere
  \end{steps}
\end{proof}

\begin{example}[Every quantum measurement defines a probability measure]\label{ex:qmeasure}
  Let a projective measurement with projectors $P_m$ ($m\in\Omega$, a finite set of outcomes) be performed on the state $\ket\psi$. Show that $\PP(B)=\bra\psi\bigl(\sum_{m\in B}P_m\bigr)\ket\psi$ is a probability measure on $\Omega$.
\end{example}
\begin{solution}
  We check the two axioms of Definition~\ref{def:probmeasure}, taking $\mathcal F$ to be all subsets of $\Omega$.
  \begin{steps}
    \item \emph{Values in $[0,1]$.} The operator $P_B=\sum_{m\in B}P_m$ is a projector, since the $P_m$ are orthogonal projectors ($P_mP_{m'}=\delta_{mm'}P_m$). Hence $\PP(B)=\bra\psi P_B\ket\psi=\norm{P_B\psi}^2\in[0,\norm\psi^2]=[0,1]$.
    \item \emph{Total mass.} $P_\Omega=\sum_mP_m=\Id$ by completeness, so $\PP(\Omega)=\braket\psi\psi=1$.
    \item \emph{Additivity.} If $B_1,B_2,\dots$ are disjoint, then $P_{\cup B_i}=\sum_iP_{B_i}$, and taking $\bra\psi\cdot\ket\psi$ gives $\PP(\cup B_i)=\sum_i\PP(B_i)$.
  \end{steps}
  By Proposition~\ref{prop:probprops}, we can now say, for instance, that the probability of getting an outcome in the set $B$ is $1$ minus the probability of getting one outside it.
\end{solution}

\begin{lookahead}
  This is the Born rule in disguise. For an observable with a \emph{continuous} spectrum, such as position, the sum over $m\in B$ is replaced by a projection $P(B)$ onto the part of the spectrum lying in the set $B$, and $B\mapsto\bra\psi P(B)\ket\psi$ is again a probability measure. This
  \emph{projection-valued measure} is how the postulates are written in full generality; we do not need it in what follows, but the example shows that nothing in this chapter is foreign to Chapter~\ref{ch:postulates}.
\end{lookahead}

\subsection{Discrete and continuous distributions}

Two kinds of probability measures account for nearly everything.

\begin{itemize}
  \item \textbf{Discrete.} $\Omega$ is finite or countable, and $\PP$ is determined by the \term{probability mass function} $p(x)=\PP(\{x\})\ge0$ with $\sum_xp(x)=1$, through $\PP(A)=\sum_{x\in A}p(x)$.
    The standard examples are \emph{Bernoulli}$(p)$ ($\Omega=\{0,1\}$, $\PP(1)=p$), \emph{uniform} on $n$ points, and \emph{binomial}$(n,p)$: the number of $1$'s in $n$ independent Bernoulli trials,
    $\PP(k)=\binom nkp^k(1-p)^{n-k}$.
  \item \textbf{Continuous.} $\Omega=\R$ and $\PP$ is given by a \term{density} $f\ge0$ with $\int_{\R}f(x)\dd x=1$, through $\PP(A)=\int_Af(x)\dd x$.
    The standard examples are the \emph{uniform} density $\frac1{b-a}$ on $[a,b]$, the \emph{exponential} density $\lambda e^{-\lambda x}$ on $[0,\infty)$, and the \emph{Gaussian} (or normal) density $N(\mu,\sigma^2)$:
    \[
      f(x)=\frac{1}{\sigma\sqrt{2\pi}}\,e^{-(x-\mu)^2/(2\sigma^2)} .
    \]
\end{itemize}
Note that for a density every single point has probability $\int_{\{x\}}f=0$: only sets of positive length carry probability.

\begin{example}\label{ex:laplace}
  Find $c$ such that $f(x)=c\,e^{-\abs x}$ is a probability density on $\R$, and compute $\PP(\abs X\le1)$.
\end{example}
\begin{solution}
  \begin{steps}
    \item \emph{Normalise.} By symmetry, $\int_{\R}e^{-\abs x}\dd x=2\int_0^\infty e^{-x}\dd x=2$, so $c=\tfrac12$.
    \item \emph{Compute the probability.} $\PP(\abs X\le1)=\int_{-1}^1\tfrac12e^{-\abs x}\dd x=\int_0^1e^{-x}\dd x=1-e^{-1}$.
    \item \emph{Number.} $1-e^{-1}\approx0.632$.
  \end{steps}
\end{solution}

\begin{example}[Born rule for a particle in a box]\label{ex:box}
  A particle in a box $[0,1]$ is in the state with wave function $\psi(x)=\sqrt2\sin(\pi x)$. Its position has density $\abs{\psi(x)}^2$. Check normalisation and find the probability that the particle is found in $[\tfrac14,\tfrac34]$.
\end{example}
\begin{solution}
  \begin{steps}
    \item \emph{The density.} $\abs{\psi(x)}^2=2\sin^2(\pi x)=1-\cos(2\pi x)$.
    \item \emph{Normalisation.} $\int_0^1\bigl(1-\cos2\pi x\bigr)\dd x=1-\bigl[\tfrac{\sin2\pi x}{2\pi}\bigr]_0^1=1$.
    \item \emph{The probability.} $\int_{1/4}^{3/4}\bigl(1-\cos2\pi x\bigr)\dd x=\tfrac12-\tfrac1{2\pi}\bigl(\sin\tfrac{3\pi}2-\sin\tfrac\pi2\bigr)=\tfrac12+\tfrac1\pi\approx0.818$.
    \item \emph{Reading.} The middle half of the box carries about $82\%$ of the probability: the particle is most likely found near the centre, where $\abs\psi^2$ is largest.
  \end{steps}
\end{solution}

\subsection{Random variables, expectation and variance}

\begin{definition}[Random variable, expectation, variance]\label{def:rv}
  A \term{random variable} is a function $X:\Omega\to\R$ (measurable, which is automatic in all our examples). Its \term{distribution} is the probability measure $B\mapsto\PP(X\in B)$ on $\R$. Its \term{expectation} and \term{variance} are
  \[
    \EE[X]=\sum_xx\,p(x)\ \text{ or }\ \int_{\R}x\,f(x)\dd x,\qquad
    \Var(X)=\EE\bigl[(X-\EE X)^2\bigr]=\EE[X^2]-\EE[X]^2 ,
  \]
  in the discrete, respectively continuous, case. The \term{standard deviation} is $\sqrt{\Var(X)}$.
\end{definition}

For a function $g$, $\EE[g(X)]=\sum_xg(x)p(x)$ or $\int g(x)f(x)\dd x$. Expectation is linear, $\EE[aX+bY]=a\EE X+b\EE Y$, and $\Var(aX+b)=a^2\Var(X)$.

\begin{example}[Quantum measurement outcomes as random variables]\label{ex:rv-qm}
  Measure the observable $Z$ on a qubit with Bloch vector $\vec n$. Compute $\EE[X]$ and $\Var(X)$ for the outcome $X\in\{\pm1\}$, and compare with $\langle Z\rangle$ and $\Delta Z$ (Section~\ref{sec:measurement}).
\end{example}
\begin{solution}
  \begin{steps}
    \item \emph{The distribution.} $\PP(X=+1)=\tfrac12(1+n_z)$ and $\PP(X=-1)=\tfrac12(1-n_z)$ (Proposition~\ref{prop:bloch}).
    \item \emph{Expectation.} $\EE[X]=\tfrac12(1+n_z)-\tfrac12(1-n_z)=n_z=\langle Z\rangle$.
    \item \emph{Variance.} $X^2=1$ always, so $\Var(X)=1-n_z^2$, which is $(\Delta Z)^2$.
  \end{steps}
  The expectation value and standard deviation of Chapter~\ref{ch:postulates} are nothing but the expectation and standard deviation of the outcome, as a random variable.
\end{solution}

How concentrated is a random variable around its mean? The variance gives a quantitative answer.

\begin{theorem}[Chebyshev's inequality]\label{thm:chebyshev}
  For every random variable $X$ with finite variance and every $\varepsilon>0$,
  \[
    \PP\bigl(\abs{X-\EE X}\ge\varepsilon\bigr)\le\frac{\Var(X)}{\varepsilon^2}.
  \]
\end{theorem}
\begin{proof}
  Put $Y=(X-\EE X)^2\ge0$ and let $A$ be the event $\{Y\ge\varepsilon^2\}$.
  \begin{steps}
    \item Pointwise, $Y\ge\varepsilon^2\,\mathbf 1_A$, where $\mathbf 1_A$ is $1$ on $A$ and $0$ elsewhere.
    \item Take expectations: $\EE[Y]\ge\varepsilon^2\PP(A)$.
    \item Since $\EE[Y]=\Var(X)$ and $A=\{\abs{X-\EE X}\ge\varepsilon\}$, dividing by $\varepsilon^2$ gives the claim. \qedhere
  \end{steps}
\end{proof}

\subsection{Conditioning and independence}

\begin{definition}[Conditional probability, independence]\label{def:conditional}
  If $\PP(B)>0$, the \term{conditional probability} of $A$ given $B$ is $\PP(A\mid B)=\PP(A\cap B)/\PP(B)$. Events $A,B$ are \term{independent} if $\PP(A\cap B)=\PP(A)\PP(B)$, that is, if learning $B$ does not change the probability of $A$.
\end{definition}

Two consequences are used constantly. If $B_1,\dots,B_n$ partition $\Omega$ with $\PP(B_i)>0$, then:
\begin{steps}
  \item \emph{total probability:} $\PP(A)=\sum_i\PP(A\mid B_i)\PP(B_i)$;
  \item \emph{Bayes' rule:} $\PP(B_j\mid A)=\dfrac{\PP(A\mid B_j)\PP(B_j)}{\sum_i\PP(A\mid B_i)\PP(B_i)}$.
\end{steps}
The first follows from $A=\bigcup_i(A\cap B_i)$ and additivity; the second is the definition of conditional probability, $\PP(B_j\mid A)=\PP(A\cap B_j)/\PP(A)$, with the numerator rewritten as $\PP(A\mid B_j)\PP(B_j)$ and the denominator by total probability.

\begin{example}[Bob guesses the state]\label{ex:bayes}
  Alice sends $\ket0$ or $\ket+$, each with probability $\tfrac12$. Bob measures $Z$. For each outcome, what does he conclude?
\end{example}
\begin{solution}
  Let $S_0,S_+$ be the events ``$\ket0$ was sent'', ``$\ket+$ was sent'', and $O_0,O_1$ the events ``outcome $0$'', ``outcome $1$''.
  \begin{steps}
    \item \emph{Likelihoods.} $\PP(O_0\mid S_0)=1$, $\PP(O_1\mid S_0)=0$; $\PP(O_0\mid S_+)=\PP(O_1\mid S_+)=\tfrac12$.
    \item \emph{Total probability of each outcome.} $\PP(O_0)=1\cdot\tfrac12+\tfrac12\cdot\tfrac12=\tfrac34$, $\PP(O_1)=0+\tfrac12\cdot\tfrac12=\tfrac14$.
    \item \emph{Bayes for outcome $0$.} $\PP(S_0\mid O_0)=\dfrac{1\cdot\frac12}{3/4}=\tfrac23$ and so $\PP(S_+\mid O_0)=\tfrac13$.
    \item \emph{Bayes for outcome $1$.} $\PP(S_0\mid O_1)=\dfrac{0\cdot\frac12}{1/4}=0$, so $\PP(S_+\mid O_1)=1$.
  \end{steps}
  Outcome $1$ is conclusive: it cannot come from $\ket0$. Outcome $0$ only makes $\ket0$ more likely ($\tfrac23$ against $\tfrac12$). This is the same phenomenon that underlies the ``never wrong, sometimes silent'' measurement of Example~\ref{ex:unambiguous}.
\end{solution}

\paragraph{Joint distributions and marginals.}
A pair of random variables $(X,Y)$ has a \term{joint distribution} $p(x,y)$; its \term{marginals} are $p_X(x)=\sum_yp(x,y)$ and $p_Y(y)=\sum_xp(x,y)$, and $X,Y$ are independent iff $p(x,y)=p_X(x)p_Y(y)$ for all $x,y$. The \term{product measure} $\PP_1\otimes\PP_2$ on $\Omega_1\times\Omega_2$ is the joint distribution of two independent variables.
This is exactly the probabilistic counterpart of Postulate~4: product states produce product measures, and entangled states produce joint distributions that are \emph{not} products.

\begin{example}\label{ex:joint-bell}
  Measure $Z$ on both qubits of (a) $\ket+\ket+$ and (b) $\ket{\Phi^+}=\frac1{\sqrt2}(\ket{00}+\ket{11})$. Write the joint distribution and decide whether the outcomes are independent.
\end{example}
\begin{solution}
  The joint probability of outcomes $(a,b)$ is $p(a,b)=\abs{\braket{ab}{\psi}}^2$.
  \begin{steps}
    \item \emph{(a)} $\ket+\ket+=\tfrac12(\ket{00}+\ket{01}+\ket{10}+\ket{11})$, so $p(a,b)=\tfrac14$ for all four pairs. The marginals are $\tfrac12,\tfrac12$ and $\tfrac14=\tfrac12\cdot\tfrac12$: the outcomes are independent.
    \item \emph{(b)} $p(0,0)=p(1,1)=\tfrac12$ and $p(0,1)=p(1,0)=0$. The marginals are again $\tfrac12,\tfrac12$, but $p(0,0)=\tfrac12\ne\tfrac14$: the outcomes are perfectly correlated, not independent.
  \end{steps}
\end{solution}

\begin{pitfall}
  The table in part (b) is exactly the one produced by a \emph{classical} mixture: a fair coin decides whether the pair is $\ket{00}$ or $\ket{11}$. Correlations in a single basis cannot prove entanglement. What distinguishes the two is the behaviour in \emph{other} bases, which is the content of Bell's inequality (Section~\ref{sec:bell}).
\end{pitfall}

\subsection{Repeated experiments: the law of large numbers}

How do we measure a probability in the laboratory? We repeat the experiment $n$ times and count. Let $X_1,\dots,X_n$ be independent Bernoulli$(p)$ variables (the $i$th is $1$ if the outcome of interest occurred in run $i$), and $\hat p=\frac1n\sum_iX_i$ the observed frequency.

\begin{theorem}[Weak law of large numbers, Bernoulli case]\label{thm:wlln}
  For every $\varepsilon>0$, $\ \PP\bigl(\abs{\hat p-p}\ge\varepsilon\bigr)\le\dfrac{1}{4n\varepsilon^2}$. In particular $\hat p\to p$ in probability as $n\to\infty$.
\end{theorem}
\begin{proof}
  \begin{steps}
    \item $\EE[\hat p]=\frac1n\sum_i\EE[X_i]=p$, by linearity.
    \item Independence makes variances add: $\Var(\hat p)=\frac1{n^2}\sum_i\Var(X_i)=\frac{p(1-p)}n$.
    \item Since $p(1-p)\le\tfrac14$, Chebyshev's inequality gives $\PP(\abs{\hat p-p}\ge\varepsilon)\le\frac{p(1-p)}{n\varepsilon^2}\le\frac1{4n\varepsilon^2}$. \qedhere
  \end{steps}
\end{proof}

\begin{example}\label{ex:sample-size}
  How many runs are enough to estimate an outcome probability within $\pm0.05$ with confidence at least $95\%$?
\end{example}
\begin{solution}
  We need $\frac1{4n\varepsilon^2}\le0.05$ with $\varepsilon=0.05$. Then $n\ge\frac1{4\cdot0.0025\cdot0.05}=2000$. So about two thousand runs suffice (Chebyshev is conservative, so fewer would do in practice).
\end{solution}

\begin{lookahead}
  This is how quantum states are reconstructed in practice (\emph{tomography}). To find the Bloch vector of an unknown qubit, we prepare many copies and measure $X$, $Y$ and $Z$ on different batches; the frequencies estimate $n_x,n_y,n_z=\langle\sigma_j\rangle$, and the error shrinks like $1/\sqrt n$. It also explains why a quantum algorithm which gives the right answer with probability $\tfrac23$ can be made reliable by simple repetition.
\end{lookahead}

\begin{exerciseset}
  \exercise{1} Show that the power set of a finite set is a $\sigma$-algebra. Show that if $A,B\in\mathcal F$ then $A\cap B$ and $A\setminus B$ are in $\mathcal F$.

  \exercise{1} Prove parts (a)--(e) of Proposition~\ref{prop:probprops} for a probability measure on a finite set, using only finite additivity and $\PP(\Omega)=1$.

  \exercise{2} Prove the inclusion--exclusion formula for three events, $\PP(A\cup B\cup C)=\PP(A)+\PP(B)+\PP(C)-\PP(A\cap B)-\PP(A\cap C)-\PP(B\cap C)+\PP(A\cap B\cap C)$.

  \exercise{2} Show that $B\mapsto\bra\psi P_B\ket\psi$ is a probability measure when $\{P_m\}$ is a complete family of orthogonal projectors (Example~\ref{ex:qmeasure}), and compute it for the measurement of $X$ on $\ket\psi=\frac1{\sqrt3}\ket0+\frac{1+\ii}{\sqrt3}\ket1$.

  \exercise{1} Find $c$ so that $f(x)=cx^2$ is a probability density on $[0,1]$, and compute $\EE[X]$ and $\Var(X)$.

  \exercise{2} Let $X$ be a Gaussian $N(\mu,\sigma^2)$. Show that $\EE[X]=\mu$ and $\Var(X)=\sigma^2$ (use $\int e^{-t^2/2}\dd t=\sqrt{2\pi}$ and integrate by parts).

  \exercise{2} A particle in a box has wave function $\psi(x)=\sqrt2\sin(2\pi x)$ on $[0,1]$. Check normalisation and compute the probability of finding it in $[0,\tfrac14]$ and the expectation of its position.

  \exercise{2} For a projective measurement of $M=\sum_mmP_m$ in the state $\ket\psi$, show that the outcome $X$ satisfies $\EE[X]=\bra\psi M\ket\psi$, $\EE[f(X)]=\bra\psi f(M)\ket\psi$ and $\Var(X)=(\Delta M)^2$.

  \exercise{2} Prove Markov's inequality, $\PP(X\ge a)\le\EE[X]/a$ for $X\ge0$ and $a>0$, and deduce Chebyshev's inequality from it.

  \exercise{2} Let $\ket\psi=\frac1{\sqrt3}\ket0+\sqrt{\tfrac23}\ket1$ be sent with probability $\tfrac13$ and $\ket{-}$ with probability $\tfrac23$. Bob measures $X$. Find the posterior probability of each state given each outcome.

  \exercise{2} Show that $A,B$ independent implies $A,B^c$ independent. Compute the joint distribution of $Z\otimes Z$ outcomes for $\ket{\Phi^+}$, $\bell01$, and for $\ket0\ket+$, and decide independence in each case.

  \exercise{3} A classical mixture $\rho=\tfrac12\ket{00}\!\bra{00}+\tfrac12\ket{11}\!\bra{11}$ and the state $\ket{\Phi^+}$ give the same joint distribution for $Z\otimes Z$. Compute the joint distributions of outcomes of $X\otimes X$ for both, and conclude that $X\otimes X$ distinguishes them.

  \exercise{2} A qubit is measured $n=100$ times in the computational basis and $\ket1$ occurs $37$ times. Estimate $p(1)$ and give a Chebyshev bound on the probability that your estimate is off by more than $0.1$.

  \exercise{3} Show that the sample mean of $n$ independent Bernoulli$(p)$ variables has the binomial law $\PP(\hat p=k/n)=\binom nkp^k(1-p)^{n-k}$, and verify that $\EE[\hat p]=p$ and $\Var(\hat p)=p(1-p)/n$ directly from the binomial sums.
\end{exerciseset}

%======================================================================
\section{Fourier transforms}\label{sec:fourier}
%======================================================================

\begin{objectives}
  You will expand periodic functions in Fourier series, compute Fourier transforms and use their main properties; relate the width of a function to the width of its transform;
  use the discrete Fourier transform as a unitary matrix; and see the characteristic function of a probability measure as a Fourier transform.
\end{objectives}

\begin{recall}
  The inner product $\langle f,g\rangle=\int\overline fg$, orthonormal bases and Parseval's identity (Section~\ref{sec:hilbert}); unitary operators (Section~\ref{sec:operators}); the momentum operator (Section~\ref{sec:operators}, Remark) and the uncertainty relation (Section~\ref{sec:measurement}); unitary representations (Section~\ref{sec:grouprep}); $\omega=e^{2\pi\ii/N}$ and roots of unity (Section~\ref{sec:complex}).
\end{recall}

A wave function $\psi(x)$ tells us where a particle is likely to be found. The \emph{same} information can be encoded in a second function $\hat\psi(k)$, which tells us how much of each wavelength the wave contains, and therefore, as we shall see, what momentum
the particle is likely to have. The dictionary between $\psi$ and $\hat\psi$ is the Fourier transform. It is a change of orthonormal basis in an infinite-dimensional Hilbert space, so the language of Chapter~\ref{ch:tools} applies directly.

\subsection{Fourier series}

Consider the space $L^2([0,1])$ of (complex) functions $f$ on $[0,1]$ with $\int_0^1\abs{f}^2<\infty$, with the inner product $\langle f,g\rangle=\int_0^1\overline{f(x)}\,g(x)\dd x$ of Section~\ref{sec:hilbert}. For $n\in\Z$ put $e_n(x)=e^{2\pi\ii nx}$.

\begin{proposition}[An orthonormal system]\label{prop:expo-orthonormal}
  The functions $e_n$, $n\in\Z$, are orthonormal: $\langle e_m,e_n\rangle=\delta_{mn}$.
\end{proposition}
\begin{proof}
  \begin{steps}
    \item $\langle e_m,e_n\rangle=\int_0^1e^{-2\pi\ii mx}e^{2\pi\ii nx}\dd x=\int_0^1e^{2\pi\ii(n-m)x}\dd x$.
    \item If $n=m$ the integrand is $1$, and the integral is $1$.
    \item If $n\neq m$, the integral is $\bigl[\frac{e^{2\pi\ii(n-m)x}}{2\pi\ii(n-m)}\bigr]_0^1=\frac{1-1}{2\pi\ii(n-m)}=0$. \qedhere
  \end{steps}
\end{proof}

It is a fact (not proved here) that this system is also \emph{complete}: it is an orthonormal basis of $L^2([0,1])$. By Theorem~\ref{thm:parseval} (valid also in infinite dimension) every $f\in L^2$ is then expanded, and Parseval's identity holds:
\[
  f=\sum_{n\in\Z}c_ne_n,\quad c_n=\langle e_n,f\rangle=\int_0^1f(x)e^{-2\pi\ii nx}\dd x,\qquad \int_0^1\abs{f}^2\dd x=\sum_{n\in\Z}\abs{c_n}^2 .
\]
This is a \term{Fourier series}: $f$ is decomposed into pure oscillations $e^{2\pi\ii nx}$ of frequency $n$, with amplitudes $c_n$.

\begin{example}[Fourier coefficients of a ramp, and a surprise]\label{ex:fourier-ramp}
  Compute the Fourier coefficients of $f(x)=x$ on $[0,1)$, and use Parseval's identity to evaluate $\sum_{n\ge1}1/n^2$.
\end{example}
\begin{solution}
  \begin{steps}
    \item \emph{The zeroth coefficient.} $c_0=\int_0^1x\dd x=\tfrac12$.
    \item \emph{The others ($n\neq0$).} Integrate by parts with $u=x$, $\dd v=e^{-2\pi\ii nx}\dd x$:
      \[
        c_n=\Bigl[\frac{x\,e^{-2\pi\ii nx}}{-2\pi\ii n}\Bigr]_0^1+\frac1{2\pi\ii n}\int_0^1e^{-2\pi\ii nx}\dd x
        =\frac{1}{-2\pi\ii n}+0=\frac{\ii}{2\pi n},
      \]
      since $e^{-2\pi\ii n}=1$ and the last integral vanishes.
    \item \emph{The norm.} $\int_0^1x^2\dd x=\tfrac13$.
    \item \emph{Parseval.} $\tfrac13=\abs{c_0}^2+\sum_{n\neq0}\abs{c_n}^2=\tfrac14+2\sum_{n\ge1}\frac1{4\pi^2n^2}$.
    \item \emph{Solve.} Hence $\tfrac{1}{2\pi^2}\sum_{n\ge1}n^{-2}=\tfrac13-\tfrac14=\tfrac1{12}$, that is,
      \[
        \sum_{n=1}^\infty\frac1{n^2}=\frac{\pi^2}{6}.
      \]
  \end{steps}
  A statement about orthonormal bases in a Hilbert space has produced a classical number-theoretic identity.
\end{solution}

\subsection{The Fourier transform on $\R$}

For functions on the whole line there is no period, and the discrete frequencies $n$ become a continuous frequency (wave number) $k\in\R$. We use the symmetric, unitary normalisation.

\begin{definition}[Fourier transform]\label{def:ft}
  For $f\in L^1(\R)\cap L^2(\R)$, its \term{Fourier transform} is
  \[
    \hat f(k)=(\FT f)(k)=\frac1{\sqrt{2\pi}}\int_{\R}f(x)\,e^{-\ii kx}\dd x,
    \qquad\text{and}\qquad
    f(x)=\frac1{\sqrt{2\pi}}\int_{\R}\hat f(k)\,e^{\ii kx}\dd k
  \]
  is the \term{inversion formula}.
\end{definition}

The function $\hat f(k)$ is the amplitude of the pure wave $e^{\ii kx}$ inside $f$, in the same way as $c_n$ above; the inversion formula is the continuous analogue of ``$f=\sum c_ne_n$''. The most important theorem is the analogue of Parseval's identity.

\begin{theorem}[Plancherel]\label{thm:plancherel}
  $\FT$ extends to a \emph{unitary} operator on $L^2(\R)$: $\ \norm{\hat f}=\norm f$ and $\langle\hat f,\hat g\rangle=\langle f,g\rangle$.
\end{theorem}

We use this without proof. Unitary is the key word: \emph{the Fourier transform is just a change of orthonormal basis}, from the ``position'' basis to the ``wave-number'' basis, and, as for any unitary, it preserves norms, and therefore probabilities.
The properties that make it useful turn calculus into algebra.

\begin{proposition}[Properties of the Fourier transform]\label{prop:ftprops}
  Let $f,g$ be sufficiently regular, with $f$ and $f'$ vanishing at $\pm\infty$. Then
  \begin{enumerate}
    \item[(a)] \emph{linearity:} $\widehat{af+bg}=a\hat f+b\hat g$;
    \item[(b)] \emph{translation:} if $(\tau_af)(x)=f(x-a)$, then $\widehat{\tau_af}(k)=e^{-\ii ka}\hat f(k)$;
    \item[(c)] \emph{derivative:} $\widehat{f'}(k)=\ii k\,\hat f(k)$;
    \item[(d)] \emph{multiplication by $x$:} $\widehat{xf}(k)=\ii\,\hat f\,'(k)$;
    \item[(e)] \emph{scaling:} for $a\neq0$, $\widehat{f(a\,\cdot)}(k)=\frac1{\abs a}\hat f(k/a)$.
  \end{enumerate}
\end{proposition}
\begin{proof}
  \begin{steps}
    \item (a) The integral is linear.
    \item (b) Substitute $y=x-a$: $\frac1{\sqrt{2\pi}}\int f(x-a)e^{-\ii kx}\dd x=\frac1{\sqrt{2\pi}}\int f(y)e^{-\ii k(y+a)}\dd y=e^{-\ii ka}\hat f(k)$.
    \item (c) Integrate by parts: $\int f'(x)e^{-\ii kx}\dd x=\bigl[f(x)e^{-\ii kx}\bigr]_{-\infty}^\infty+\ii k\int f(x)e^{-\ii kx}\dd x$, and the boundary term vanishes.
    \item (d) Differentiate under the integral sign: $\hat f\,'(k)=\frac1{\sqrt{2\pi}}\int f(x)(-\ii x)e^{-\ii kx}\dd x=-\ii\,\widehat{xf}(k)$, and multiply by $\ii$.
    \item (e) Substitute $y=ax$ (for $a>0$; the case $a<0$ is the same up to orientation): $\frac1{\sqrt{2\pi}}\int f(ax)e^{-\ii kx}\dd x=\frac1{a}\frac1{\sqrt{2\pi}}\int f(y)e^{-\ii(k/a)y}\dd y$. \qedhere
  \end{steps}
\end{proof}

Property (c) is the one that matters most to physics: \emph{differentiation becomes multiplication}. This is why the Fourier transform solves linear differential equations with constant coefficients, such as Schr\"odinger's free equation.

\paragraph{The Gaussian is its own transform.}
Gaussians are the one family that Fourier transforms map to themselves, and they will turn out to be the states of minimal uncertainty.

\begin{example}\label{ex:gaussian-ft}
  Show that $f(x)=e^{-x^2/2}$ has $\hat f(k)=e^{-k^2/2}$, and deduce the transform of $g(x)=e^{-x^2/(2\sigma^2)}$.
\end{example}
\begin{solution}
  \begin{steps}
    \item \emph{A differential equation for $f$.} $f'(x)=-x\,f(x)$.
    \item \emph{Transform it.} Apply $\FT$, with (c) on the left and (d) on the right: $\ii k\hat f(k)=-\ii\,\hat f\,'(k)$, that is, $\hat f\,'(k)=-k\,\hat f(k)$.
    \item \emph{Solve.} The solutions are $\hat f(k)=\hat f(0)\,e^{-k^2/2}$.
    \item \emph{The constant.} $\hat f(0)=\frac1{\sqrt{2\pi}}\int e^{-x^2/2}\dd x=\frac{\sqrt{2\pi}}{\sqrt{2\pi}}=1$. So $\hat f(k)=e^{-k^2/2}$.
    \item \emph{Rescale.} $g(x)=f(x/\sigma)$, which is $f(a\,\cdot)$ with $a=1/\sigma$. By (e), $\hat g(k)=\sigma\,\hat f(\sigma k)=\sigma\,e^{-\sigma^2k^2/2}$.
  \end{steps}
  Note the reciprocal relation: a Gaussian of width $\sigma$ in $x$ becomes a Gaussian of width $1/\sigma$ in $k$.
\end{solution}

\subsection{Width in $x$ against width in $k$}

Property (e) says it all: compressing a function by $a>1$ stretches its transform by the same factor. A narrow function needs many wavelengths to be built; a wide one needs few. This trade-off is quantitative.

\begin{example}[The Gaussian wave packet saturates the uncertainty relation]\label{ex:gaussian-uncertainty}
  Let $\psi(x)=(\pi\sigma^2)^{-1/4}e^{-x^2/(2\sigma^2)}$. Check that $\psi$ and $\hat\psi$ are normalised, compute the standard deviations $\Delta x$ of $\abs\psi^2$ and $\Delta k$ of $\abs{\hat\psi}^2$, and show $\Delta x\,\Delta k=\tfrac12$.
\end{example}
\begin{solution}
  \begin{steps}
    \item \emph{The position density.} $\abs{\psi(x)}^2=\frac1{\sigma\sqrt\pi}e^{-x^2/\sigma^2}$, a Gaussian $N(0,s^2)$ with $2s^2=\sigma^2$. It integrates to $1$, has mean $0$ and variance $\sigma^2/2$, so $\Delta x=\sigma/\sqrt2$.
    \item \emph{The transform.} By Example~\ref{ex:gaussian-ft}, $\hat\psi(k)=(\pi\sigma^2)^{-1/4}\sigma\,e^{-\sigma^2k^2/2}=(\sigma^2/\pi)^{1/4}e^{-\sigma^2k^2/2}$.
    \item \emph{The wave-number density.} $\abs{\hat\psi(k)}^2=\frac\sigma{\sqrt\pi}e^{-\sigma^2k^2}$, a Gaussian with $2s_k^2=1/\sigma^2$: it integrates to $1$ (as Plancherel promises), has mean $0$ and variance $\frac1{2\sigma^2}$, so $\Delta k=\frac1{\sigma\sqrt2}$.
    \item \emph{The product.} $\Delta x\,\Delta k=\frac\sigma{\sqrt2}\cdot\frac1{\sigma\sqrt2}=\frac12$, whatever the value of $\sigma$.
  \end{steps}
\end{solution}

The momentum of a particle is $p=\hbar k$, so $\Delta x\,\Delta p=\hbar/2$: the Gaussian attains the lower bound of the Heisenberg--Robertson relation (Theorem~\ref{thm:uncertainty}), which for $[\hat x,\hat p]=\ii\hbar$ reads $\Delta x\,\Delta p\ge\hbar/2$. For other shapes the product is larger. The uncertainty principle is, in this light, a theorem about Fourier transforms.

\begin{lookahead}
  Since $\FT$ is unitary and $\widehat{f'}=\ii k\hat f$, the momentum operator $\hat p=-\ii\hbar\frac{d}{dx}$ becomes \emph{multiplication by $\hbar k$} in the wave-number representation. So the kinetic energy $\hat p^2/2m$ is the multiplication operator $\hbar^2k^2/2m$, and
  the free-particle propagator $e^{-\ii\hat Ht/\hbar}$ (Proposition~\ref{prop:schrodinger}) multiplies $\hat\psi(k)$ by the phase $e^{-\ii\hbar k^2t/2m}$. The free Schr\"odinger equation of Section~\ref{sec:motivation} is solved by one Fourier transform, one multiplication by a phase, and one inverse transform.
  This is the spectral theorem for $\hat p$, in action.
\end{lookahead}

\subsection{The discrete Fourier transform}

On $\C^N$ the Fourier transform becomes a finite matrix. Let $\omega=e^{2\pi\ii/N}$.

\begin{definition}[Discrete Fourier transform]\label{def:dft}
  The \term{discrete Fourier transform} is the $N\times N$ matrix $F_N$ with entries $(F_N)_{jk}=\frac1{\sqrt N}\,\omega^{jk}$ for $j,k=0,\dots,N-1$.
\end{definition}

\begin{theorem}\label{thm:dft-unitary}
  $F_N$ is unitary.
\end{theorem}
\begin{proof}
  We show that the columns of $F_N$ are orthonormal.
  \begin{steps}
    \item The inner product of columns $k$ and $k'$ is $\sum_j\overline{\tfrac{\omega^{jk}}{\sqrt N}}\tfrac{\omega^{jk'}}{\sqrt N}=\frac1N\sum_{j=0}^{N-1}\omega^{j(k'-k)}$.
    \item If $k=k'$ every term is $1$, so the sum is $N$ and the inner product is $1$.
    \item If $k\neq k'$ put $r=\omega^{k'-k}\neq1$; since $r^N=1$, the geometric sum is $\sum_jr^j=\frac{1-r^N}{1-r}=0$. \qedhere
  \end{steps}
\end{proof}

\begin{example}\label{ex:dft}
  Write $F_2$ and $F_4$, and compute $F_4$ applied to $\frac12(1,1,1,1)^t$ and to $(1,0,0,0)^t$.
\end{example}
\begin{solution}
  \begin{steps}
    \item \emph{$F_2$.} $\omega=e^{\ii\pi}=-1$, so $F_2=\frac1{\sqrt2}\begin{pmatrix}1&1\\1&-1\end{pmatrix}$: \emph{the Hadamard gate}.
    \item \emph{$F_4$.} Here $\omega=\ii$, and the entries are $\ii^{jk}/2$:
      \[
        F_4=\frac12\begin{pmatrix}1&1&1&1\\1&\ii&-1&-\ii\\1&-1&1&-1\\1&-\ii&-1&\ii\end{pmatrix}.
      \]
    \item \emph{Uniform vector.} $F_4\cdot\frac12(1,1,1,1)^t$: the first row gives $\frac14\cdot4=1$; every other row sums to $0$ (the sum of the fourth roots of unity vanishes). The result is $(1,0,0,0)^t$.
    \item \emph{Basis vector.} $F_4(1,0,0,0)^t$ is the first column, $\tfrac12(1,1,1,1)^t$.
  \end{steps}
  So $F_4$ exchanges ``perfectly localised'' and ``perfectly spread out'' vectors: the discrete analogue of the position--momentum duality.
\end{solution}

\begin{lookahead}
  As a unitary matrix, $F_{2^n}$ is a legitimate quantum gate on $n$ qubits: the \emph{quantum Fourier transform}. For $n=1$ it is the Hadamard gate, and for general $n$ it can be implemented with about $n^2$ elementary gates, exponentially fewer than the $2^n$ entries of the matrix suggest. It is the engine of Shor's algorithm: it turns a hidden \emph{period} of a function, which is spread
  over the amplitudes, into a sharply peaked distribution that a measurement can read.
\end{lookahead}

\subsection{The characteristic function of a probability measure}

If $X$ has a density $f$, then the Fourier transform of the \emph{density} is, up to normalisation, a very useful object.

\begin{definition}[Characteristic function]\label{def:charfn}
  The \term{characteristic function} of a random variable $X$ is $\varphi_X(k)=\EE\bigl[e^{\ii kX}\bigr]$. If $X$ has density $f$, then $\varphi_X(k)=\int e^{\ii kx}f(x)\dd x=\sqrt{2\pi}\,\hat f(-k)$; if $X$ is discrete, $\varphi_X(k)=\sum_xp(x)e^{\ii kx}$.
\end{definition}

\begin{proposition}\label{prop:charfn}
  For every random variable: (a) $\varphi_X(0)=1$ and $\abs{\varphi_X(k)}\le1$; (b) $\varphi_{aX+b}(k)=e^{\ii kb}\varphi_X(ak)$; (c) if $X,Y$ are independent then $\varphi_{X+Y}=\varphi_X\varphi_Y$; (d) $\varphi_X'(0)=\ii\,\EE[X]$ when $\EE\abs X<\infty$.
\end{proposition}
\begin{proof}
  \begin{steps}
    \item (a) $e^{\ii0}=1$, and $\abs{\EE[e^{\ii kX}]}\le\EE\abs{e^{\ii kX}}=1$.
    \item (b) $\EE[e^{\ii k(aX+b)}]=e^{\ii kb}\EE[e^{\ii(ka)X}]$.
    \item (c) By independence, $\EE[e^{\ii kX}e^{\ii kY}]=\EE[e^{\ii kX}]\,\EE[e^{\ii kY}]$.
    \item (d) Differentiate $\EE[e^{\ii kX}]$ at $k=0$: $\EE[\ii Xe^{\ii kX}]\big|_{k=0}=\ii\,\EE X$. \qedhere
  \end{steps}
\end{proof}

\begin{example}\label{ex:charfn}
  Find $\varphi$ for a Bernoulli$(p)$ variable and for $N(0,\sigma^2)$, and show that the sum of two independent $N(0,\sigma_1^2)$ and $N(0,\sigma_2^2)$ variables is $N(0,\sigma_1^2+\sigma_2^2)$.
\end{example}
\begin{solution}
  \begin{steps}
    \item \emph{Bernoulli.} $\varphi(k)=(1-p)e^{0}+p\,e^{\ii k}=1-p+pe^{\ii k}$.
    \item \emph{Gaussian.} The density is $f(x)=\frac1{\sigma\sqrt{2\pi}}e^{-x^2/(2\sigma^2)}$, so $\varphi(k)=\int e^{\ii kx}f(x)\dd x=\sqrt{2\pi}\hat f(-k)$. By Example~\ref{ex:gaussian-ft}, $\hat f(k)=\frac1{\sigma\sqrt{2\pi}}\sigma e^{-\sigma^2k^2/2}$, giving $\varphi(k)=e^{-\sigma^2k^2/2}$.
    \item \emph{The sum.} By (c), $\varphi_{X+Y}(k)=e^{-\sigma_1^2k^2/2}e^{-\sigma_2^2k^2/2}=e^{-(\sigma_1^2+\sigma_2^2)k^2/2}$, the characteristic function of $N(0,\sigma_1^2+\sigma_2^2)$. A characteristic function determines the distribution (a standard theorem, the injectivity of the Fourier transform), which proves the claim.
  \end{steps}
\end{solution}

A beautiful connection with quantum mechanics: for an observable $M=\sum_mmP_m$ measured in the state $\ket\psi$ with outcome $X$,
\[
  \varphi_X(k)=\sum_mp(m)e^{\ii km}=\bra\psi e^{\ii kM}\ket\psi .
\]
The characteristic function of the measurement statistics is the expectation value of the \emph{unitary} $e^{\ii kM}$. For $M=Z$, for example, $\bra\psi e^{\ii kZ}\ket\psi=\cos k+\ii\,n_z\sin k$, the characteristic function of the $\pm1$ variable with mean $n_z$.

\begin{lookahead}
  The map $k\mapsto e^{\ii kM}$ is a \emph{unitary representation of the group} $(\R,+)$ (Section~\ref{sec:grouprep}), and the statement above says that all the statistics of $M$ are encoded in the expectation of this representation. This is the viewpoint of Stone's theorem, and the starting point of the
  Fourier-analytic approach to quantum measurement.
\end{lookahead}

\begin{exerciseset}
  \exercise{1} Check that $\{e_n\}$ is an orthonormal system (Proposition~\ref{prop:expo-orthonormal}) by computing $\langle e_1,e_1\rangle$, $\langle e_1,e_2\rangle$ and $\langle e_0,e_{-3}\rangle$.

  \exercise{2} Compute the Fourier coefficients of $f(x)=\cos(2\pi x)$ and of $f(x)=\sin^2(\pi x)$ on $[0,1]$ (use $e^{\ii\theta}$ to rewrite them), and verify Parseval's identity.

  \exercise{2} Compute the Fourier coefficients of $f(x)=\mathbf 1_{[0,1/2)}(x)$ (the square wave). Which coefficients vanish? Use Parseval to evaluate $\sum_{n\ \mathrm{odd}}1/n^2$.

  \exercise{2} Using Example~\ref{ex:fourier-ramp}, show that $\sum_{n\ge1}\frac{\sin(2\pi nx)}{n}=\frac{\pi}{2}(1-2x)$ for $0<x<1$ is the (formal) Fourier series of the ramp, by combining $c_n$ and $c_{-n}$.

  \exercise{2} Prove parts (b) and (e) of Proposition~\ref{prop:ftprops}, and verify them for $f(x)=e^{-x^2/2}$.

  \exercise{2} Compute the Fourier transform of $f(x)=\mathbf 1_{[-1,1]}(x)$ and show that it is $\sqrt{2/\pi}\,\frac{\sin k}{k}$. Where is $\hat f$ concentrated when the interval is narrowed to $[-a,a]$?

  \exercise{2} Show that if $f$ is real and even, then $\hat f$ is real and even.

  \exercise{3} The \emph{convolution} is $(f*g)(x)=\int f(y)g(x-y)\dd y$. Show that $\widehat{f*g}=\sqrt{2\pi}\,\hat f\,\hat g$, and use it to show that the convolution of two Gaussians $N(0,\sigma_1^2)$ and $N(0,\sigma_2^2)$ is $N(0,\sigma_1^2+\sigma_2^2)$.

  \exercise{2} For the Gaussian packet of Example~\ref{ex:gaussian-uncertainty}, verify all the integrals, and compute $\Delta x\,\Delta p$ with $p=\hbar k$. Repeat for $\psi(x)=(\pi\sigma^2)^{-1/4}e^{-x^2/(2\sigma^2)}e^{\ii k_0x}$ and explain the role of $k_0$.

  \exercise{3} Free evolution: let $\hat\psi(k,0)$ be the initial transform. Show that $\hat\psi(k,t)=\hat\psi(k,0)e^{-\ii\hbar k^2t/2m}$ solves the free Schr\"odinger equation, and that $\abs{\hat\psi(k,t)}$ does not depend on $t$. Deduce that the momentum distribution is conserved.

  \exercise{1} Compute $F_3$ explicitly (with $\omega=e^{2\pi\ii/3}$) and check that its first two columns are orthonormal.

  \exercise{2} Prove that $F_N^2$ is the permutation $e_j\mapsto e_{-j\bmod N}$, and deduce $F_N^4=\Id$ and that the eigenvalues of $F_N$ lie in $\{1,-1,\ii,-\ii\}$.

  \exercise{2} Compute $F_4$ applied to $(1,0,1,0)^t$ and to $(1,1,0,0)^t$. Explain the first result in terms of periodicity.

  \exercise{2} Show that the characteristic function of $N(\mu,\sigma^2)$ is $e^{\ii\mu k-\sigma^2k^2/2}$, and use $\varphi'(0)=\ii\EE X$ and $\varphi''(0)=-\EE X^2$ to recover its mean and variance.

  \exercise{2} Let $\ket\psi$ have Bloch vector $\vec n$. Show that $\bra\psi e^{\ii k\,\vec v\cdot\vec\sigma}\ket\psi=\cos k+\ii\sin k\,(\vec n\cdot\vec v)$ for a unit vector $\vec v$, and identify the distribution that it characterises.
\end{exerciseset}

%======================================================================
\section{Shannon entropy}\label{sec:shannon}
%======================================================================

\begin{objectives}
  You will compute the Shannon entropy of a distribution and interpret it as an average surprise and as a compression rate; prove its basic properties; and work with joint and conditional entropies and the mutual information.
\end{objectives}

\begin{recall}
  Probability distributions, random variables, expectation, independence and conditional probability (Section~\ref{sec:prob}).
\end{recall}

How much information is there in the outcome of a random experiment? The answer depends only on the probabilities. A coin that always lands heads tells us nothing; a fair coin tells us one bit; a fair die tells us more. Claude Shannon~\cite{shannon1948} found the function that quantifies this, and
it turned out to have an operational meaning: it is the number of bits per symbol needed to store a stream of such outcomes.

\paragraph{Conventions.}
Logarithms are in base $2$, so entropy is measured in \emph{bits}, unless we write $\ln$. We use $0\log0=0$, which is the limit of $x\log x$ as $x\to0^+$. Distributions are finite.

\subsection{Information content and entropy}

If an event of probability $p$ occurs, how surprised should we be? A reasonable measure should be large for rare events and $0$ for certain ones, and surprises of \emph{independent} events should add (the surprise of two independent events is the sum). Since probabilities of independent events multiply, the logarithm is the natural choice.

\begin{definition}[Surprise and Shannon entropy]\label{def:shannon}
  The \term{information content} (or surprise) of an event of probability $p>0$ is $-\log p$. The \term{Shannon entropy} of a random variable $X$ with distribution $p(x)$ is its average surprise,
  \[
    \Hs(X)=\Hs(p)=-\sum_xp(x)\log p(x)=\EE\bigl[-\log p(X)\bigr].
  \]
\end{definition}

\begin{example}\label{ex:shannon-values}
  Compute the entropy of: (a) a fair coin; (b) a coin with $\PP(1)=0.1$; (c) a uniform variable on $8$ outcomes; (d) a deterministic variable.
\end{example}
\begin{solution}
  \begin{steps}
    \item (a) $-\tfrac12\log\tfrac12-\tfrac12\log\tfrac12=\tfrac12+\tfrac12=1$ bit.
    \item (b) $-0.1\log0.1-0.9\log0.9=0.1\times3.322+0.9\times0.152\approx0.332+0.137=0.469$ bits.
    \item (c) Each $p=\tfrac18$, so $\Hs=-\log\tfrac18=\log8=3$ bits.
    \item (d) A single outcome with $p=1$ contributes $-1\cdot\log1=0$: no surprise at all.
  \end{steps}
  The biased coin carries less than a bit: we can usually guess the outcome.
\end{solution}

For two outcomes with probabilities $p,1-p$ the entropy is the \term{binary entropy} $h(p)=-p\log p-(1-p)\log(1-p)$, plotted in Figure~\ref{fig:binary-entropy}. It is $0$ at $p=0,1$, symmetric about $p=\tfrac12$, and maximal ($=1$) at $p=\tfrac12$.

\begin{figure}[htb]
  \centering
  \begin{tikzpicture}
    \begin{axis}[width=0.7\linewidth, height=5.4cm, xmin=0, xmax=1, ymin=0, ymax=1.25,
        xlabel={$p$}, ylabel={$h(p)$ (bits)}, axis lines=left, clip=false,
        xtick={0,0.25,0.5,0.75,1}, ytick={0,0.5,1},
        label style={font=\footnotesize}, tick label style={font=\footnotesize}]
      \addplot[navy, very thick, samples=200, domain=0.001:0.999] {-x*ln(x)/ln(2)-(1-x)*ln(1-x)/ln(2)};
      \addplot[only marks, mark=*, mark size=1.8pt, figacc] coordinates {(0.5,1) (0.1,0.469) (0.25,0.811)};
      \node[figacc, font=\footnotesize, anchor=south] at (axis cs:0.5,1.0) {fair coin};
      \node[figacc, font=\footnotesize, anchor=west] at (axis cs:0.12,0.42) {$p=0.1$};
    \end{axis}
  \end{tikzpicture}
  \caption{The binary entropy $h(p)$. A fair coin is the most unpredictable (one bit); a very biased one is almost predictable.}
  \label{fig:binary-entropy}
\end{figure}

\begin{proposition}[Basic properties]\label{prop:shannon-props}
  For a distribution $p$ on $n$ outcomes:
  \begin{enumerate}
    \item[(a)] $\Hs(p)\ge0$, with equality iff $p$ is deterministic;
    \item[(b)] $\Hs(p)\le\log n$, with equality iff $p$ is uniform.
  \end{enumerate}
\end{proposition}
\begin{proof}
  \begin{steps}
    \item (a) Each term $-p(x)\log p(x)$ is $\ge0$ because $0\le p(x)\le1$, and vanishes only if $p(x)\in\{0,1\}$. The sum vanishes only if every $p(x)\in\{0,1\}$, that is, if one outcome has probability $1$.
    \item (b) Restrict to the outcomes with $p(x)>0$ (say $m\le n$ of them). The logarithm is \emph{concave}, so Jensen's inequality (for a concave $f$ and weights $p(x)\ge0$ summing to $1$: $\sum_xp(x)f(a_x)\le f\bigl(\sum_xp(x)a_x\bigr)$) gives
      \[
        \Hs(p)=\sum_xp(x)\log\frac1{p(x)}\le\log\Bigl(\sum_xp(x)\cdot\frac1{p(x)}\Bigr)=\log m\le\log n .
      \]
    \item Equality in Jensen requires $1/p(x)$ to be constant on the support, and $m=n$: the distribution is uniform. \qedhere
  \end{steps}
\end{proof}

\subsection{Several variables}

For a pair $(X,Y)$ with joint distribution $p(x,y)$, we use the joint entropy $\Hs(X,Y)=-\sum_{x,y}p(x,y)\log p(x,y)$, and we define the conditional entropy as the average remaining uncertainty:
$\Hs(Y\mid X)=\sum_xp_X(x)\,\Hs(Y\mid X=x)$, where $\Hs(Y\mid X=x)$ is the entropy of the conditional distribution $p(y\mid x)=p(x,y)/p_X(x)$.

\begin{proposition}[Chain rule]\label{prop:chainrule}
  $\Hs(X,Y)=\Hs(X)+\Hs(Y\mid X)$.
\end{proposition}
\begin{proof}
  \begin{steps}
    \item Write $p(x,y)=p_X(x)\,p(y\mid x)$, so $\log p(x,y)=\log p_X(x)+\log p(y\mid x)$.
    \item Multiply by $-p(x,y)$ and sum over $x,y$.
    \item The first part gives $-\sum_xp_X(x)\log p_X(x)=\Hs(X)$ (summing $p(x,y)$ over $y$).
    \item The second part gives $-\sum_xp_X(x)\sum_yp(y\mid x)\log p(y\mid x)=\sum_xp_X(x)\Hs(Y\mid X=x)=\Hs(Y\mid X)$. \qedhere
  \end{steps}
\end{proof}

How much does knowing $X$ tell us about $Y$? The answer is the reduction of the uncertainty about $Y$.

\begin{definition}[Mutual information]\label{def:mutualinfo}
  $I(X\!:\!Y)=\Hs(X)+\Hs(Y)-\Hs(X,Y)=\Hs(Y)-\Hs(Y\mid X)$.
\end{definition}

The two forms agree by the chain rule. We shall prove in Section~\ref{sec:relent} that $I(X\!:\!Y)\ge0$, with equality iff $X$ and $Y$ are independent; equivalently, \emph{conditioning never increases entropy}, $\Hs(Y\mid X)\le\Hs(Y)$.

\begin{example}[Correlated and independent qubit outcomes]\label{ex:mutual-info}
  Compute $\Hs(X),\Hs(Y),\Hs(X,Y)$ and $I(X\!:\!Y)$ for the $Z\otimes Z$ outcomes of (a) $\ket+\ket+$ and (b) $\ket{\Phi^+}$ (Example~\ref{ex:joint-bell}).
\end{example}
\begin{solution}
  \begin{steps}
    \item \emph{(a)} The joint distribution is uniform on $4$ outcomes, so $\Hs(X,Y)=2$; both marginals are fair, $\Hs(X)=\Hs(Y)=1$. Hence $I=1+1-2=0$: independent.
    \item \emph{(b)} The joint distribution is $(\tfrac12,0,0,\tfrac12)$, so $\Hs(X,Y)=1$; the marginals are fair, $\Hs(X)=\Hs(Y)=1$. Hence $I=1+1-1=1$ bit: knowing $X$ determines $Y$ completely.
  \end{steps}
\end{solution}

\subsection{Entropy as a compression rate}

The entropy is not just a convenient formula; it has an operational meaning.

\begin{theorem}[Shannon's source coding theorem, informally]\label{thm:source-coding}
  A source emitting independent symbols with distribution $p$ can be compressed to $\Hs(p)$ bits per symbol, asymptotically for long messages, and no better.
\end{theorem}

We do not prove this, but a small example shows the idea: give short codewords to likely symbols.

\begin{example}[A code that reaches the entropy]\label{ex:huffman}
  A source emits $a,b,c,d$ with probabilities $\tfrac12,\tfrac14,\tfrac18,\tfrac18$. Compute $\Hs$, and find a prefix code whose average length equals it.
\end{example}
\begin{solution}
  \begin{steps}
    \item \emph{Entropy.} $\Hs=\tfrac12\log2+\tfrac14\log4+2\cdot\tfrac18\log8=\tfrac12+\tfrac12+\tfrac34=1.75$ bits.
    \item \emph{A code.} Use $a\mapsto0$, $b\mapsto10$, $c\mapsto110$, $d\mapsto111$. No codeword is a prefix of another, so a stream of bits can be decoded unambiguously.
    \item \emph{Average length.} $\tfrac12\cdot1+\tfrac14\cdot2+\tfrac18\cdot3+\tfrac18\cdot3=1.75$ bits per symbol.
  \end{steps}
  A naive fixed-length code would use $2$ bits per symbol; here the code length of each symbol is exactly its surprise $-\log p(x)$, and so the average equals $\Hs$.
\end{solution}

\begin{lookahead}
  The quantum version of this theorem, due to Schumacher, says that $n$ copies of a quantum state $\rho$ can be compressed into $n\,\Ss(\rho)$ qubits, where $\Ss$ is the von Neumann entropy of Section~\ref{sec:vonneumann}. The classical capacity of a noisy channel has an equally precise answer: a binary symmetric channel that flips each bit with probability $\varepsilon$ can transmit reliably
  $1-h(\varepsilon)$ bits per use, which for $\varepsilon=0.1$ is $1-0.469\approx0.531$.
\end{lookahead}

\begin{exerciseset}
  \exercise{1} Compute the entropy of the distributions $(\tfrac12,\tfrac12)$, $(\tfrac13,\tfrac13,\tfrac13)$, $(0.9,0.1)$, $(\tfrac12,\tfrac14,\tfrac14)$ and $(1,0,0)$.

  \exercise{1} Show $h(p)=h(1-p)$ and find the maximum of $h$ (set $h'(p)=0$).

  \exercise{2} A die is loaded so that $6$ appears with probability $\tfrac12$ and the other faces with probability $\tfrac1{10}$ each. Compute its entropy and compare with $\log6$.

  \exercise{2} Prove Proposition~\ref{prop:shannon-props}(b) a second way: compare $\Hs(p)$ with $\log n$ using $\ln x\le x-1$ applied to $x=\frac1{np(x)}$.

  \exercise{2} Show that $\Hs$ is invariant under relabelling of the outcomes, and that adding outcomes of probability $0$ does not change it.

  \exercise{2} Let $X$ and $Y$ be independent. Show that $\Hs(X,Y)=\Hs(X)+\Hs(Y)$ and $I(X\!:\!Y)=0$.

  \exercise{2} For the joint distribution $p(0,0)=\tfrac12,\ p(0,1)=\tfrac14,\ p(1,1)=\tfrac14,\ p(1,0)=0$, compute $\Hs(X),\Hs(Y),\Hs(X,Y),\Hs(Y\mid X),\Hs(X\mid Y)$ and $I(X\!:\!Y)$.

  \exercise{2} Compute $I(X\!:\!Y)$ for the outcomes of $Z\otimes Z$ on $\bell01$ and on $\ket0\ket+$, and for the classical mixture $\tfrac12\ket{00}\!\bra{00}+\tfrac12\ket{11}\!\bra{11}$.

  \exercise{2} Let $X$ be a fair bit and $Y=X\oplus N$ with $N$ an independent Bernoulli$(\varepsilon)$ noise bit. Show $I(X\!:\!Y)=1-h(\varepsilon)$, and evaluate it for $\varepsilon=0.1$.

  \exercise{2} Find a prefix code for the distribution $(\tfrac12,\tfrac14,\tfrac18,\tfrac1{16},\tfrac1{16})$ with average length equal to its entropy.

  \exercise{3} Prove that $\Hs$ is a concave function of the distribution: $\Hs(\lambda p+(1-\lambda)q)\ge\lambda\Hs(p)+(1-\lambda)\Hs(q)$. (Hint: $x\mapsto-x\log x$ is concave.)

  \exercise{3} Show that for any function $g$, $\Hs(g(X))\le\Hs(X)$, with equality when $g$ is injective.
\end{exerciseset}

%======================================================================
\section{Von Neumann entropy}\label{sec:vonneumann}
%======================================================================

\begin{objectives}
  You will define the von Neumann entropy of a density operator through its spectrum and compute it for qubits and for subsystems; prove its basic properties; interpret it as the entropy of entanglement for pure bipartite states; and see how it differs from the classical entropy.
\end{objectives}

\begin{recall}
  Density operators, purity and the Bloch ball (Section~\ref{sec:density}); the partial trace (Section~\ref{sec:density}) and the Schmidt decomposition (Section~\ref{sec:schmidt}); the functional calculus $f(\rho)$ (Section~\ref{sec:operators}); the Shannon entropy and its properties (Section~\ref{sec:shannon}).
\end{recall}

A density operator $\rho$ describes a state of knowledge: how much do we \emph{not} know? Its eigenvalues $\lambda_i$ form a probability distribution (they are non-negative and add to $1$, Theorem~\ref{thm:density-char}), and the natural measure of uncertainty is the Shannon entropy of that distribution. This is the von Neumann entropy.

\subsection{Definition and first examples}

\begin{definition}[Von Neumann entropy]\label{def:vn}
  The \term{von Neumann entropy} of a density operator $\rho$ on a space of dimension $d$ is
  \[
    \Ss(\rho)=-\Tr\bigl(\rho\log\rho\bigr)=-\sum_{i=1}^d\lambda_i\log\lambda_i=\Hs(\lambda_1,\dots,\lambda_d),
  \]
  where $\lambda_i$ are the eigenvalues of $\rho$ (with $0\log0=0$), and $\rho\log\rho=f(\rho)$ is defined by the functional calculus (Definition~\ref{def:functional}) with $f(x)=x\log x$.
\end{definition}

The equality of the three expressions is the functional calculus again: if $\rho=\sum_i\lambda_i\ket{v_i}\!\bra{v_i}$ then $\rho\log\rho=\sum_i\lambda_i\log\lambda_i\ket{v_i}\!\bra{v_i}$, whose trace is $\sum_i\lambda_i\log\lambda_i$. The entropy depends only on the spectrum, not on the eigenvectors.

\begin{example}\label{ex:vn-basic}
  Compute $\Ss(\rho)$ for (a) a pure state; (b) the maximally mixed state $\Id/d$; (c) $\rho=\tfrac34\ket0\!\bra0+\tfrac14\ket1\!\bra1$.
\end{example}
\begin{solution}
  \begin{steps}
    \item (a) A pure state has eigenvalues $1,0,\dots,0$, so $\Ss=-1\log1=0$.
    \item (b) All $d$ eigenvalues equal $\tfrac1d$, so $\Ss=-d\cdot\tfrac1d\log\tfrac1d=\log d$. For a qubit, $1$ bit.
    \item (c) The eigenvalues are $\tfrac34,\tfrac14$, so $\Ss=h(\tfrac14)=-\tfrac14\log\tfrac14-\tfrac34\log\tfrac34\approx0.5+0.311=0.811$ bits.
  \end{steps}
\end{solution}

For a qubit, the Bloch ball of Section~\ref{sec:density} gives a closed formula.

\begin{proposition}\label{prop:qubit-entropy}
  If $\rho=\tfrac12(\Id+\vec r\cdot\vec\sigma)$ then $\Ss(\rho)=h\!\left(\tfrac{1+\abs{\vec r}}{2}\right)$.
\end{proposition}
\begin{proof}
  The eigenvalues of $\rho$ are $\tfrac12(1\pm\abs{\vec r})$ (Section~\ref{sec:density}), a distribution on two points with probabilities $q$ and $1-q$, $q=\tfrac12(1+\abs{\vec r})$. Its entropy is $h(q)$.
\end{proof}

So the entropy is a function of the \emph{distance from the centre} of the Bloch ball: $0$ on the surface (pure states), $1$ bit at the centre, and decreasing steadily in between (Figure~\ref{fig:entropy-ball}).

\begin{figure}[htb]
  \centering
  \begin{tikzpicture}
    \begin{axis}[width=0.7\linewidth, height=5.4cm, xmin=0, xmax=1, ymin=0, ymax=1.25,
        xlabel={$\abs{\vec r}$ (distance from the centre of the Bloch ball)}, ylabel={$\Ss(\rho)$ (bits)}, axis lines=left, clip=false,
        xtick={0,0.25,0.5,0.75,1}, ytick={0,0.5,1},
        label style={font=\footnotesize}, tick label style={font=\footnotesize}]
      \addplot[navy, very thick, samples=200, domain=0:0.9995] {-((1+x)/2)*ln((1+x)/2)/ln(2)-((1-x)/2)*ln((1-x)/2)/ln(2)};
      \addplot[only marks, mark=*, mark size=1.8pt, figacc] coordinates {(0,1) (1,0)};
      \node[figacc, font=\footnotesize, anchor=south west] at (axis cs:0.03,1.04) {maximally mixed (centre)};
      \node[figacc, font=\footnotesize, anchor=west] at (axis cs:1.03,0.06) {pure (surface)};
    \end{axis}
  \end{tikzpicture}
  \caption{The entropy of a qubit state as a function of its Bloch vector length: $\Ss=h\bigl(\frac{1+\abs{\vec r}}2\bigr)$.}
  \label{fig:entropy-ball}
\end{figure}

\begin{example}\label{ex:vn-mixture}
  Compute the entropy of the ensemble $\bigl\{\tfrac12,\ket0;\ \tfrac12,\ket+\bigr\}$ of Example~\ref{ex:mixture-bloch}, and compare it with the Shannon entropy of the ensemble probabilities.
\end{example}
\begin{solution}
  \begin{steps}
    \item \emph{The state.} $\rho=\tfrac14\begin{pmatrix}3&1\\1&1\end{pmatrix}$ with $\vec r=(\tfrac12,0,\tfrac12)$ and $\abs{\vec r}=\tfrac1{\sqrt2}$.
    \item \emph{The eigenvalues.} $\tfrac12\bigl(1\pm\tfrac1{\sqrt2}\bigr)\approx0.854,\ 0.146$.
    \item \emph{The entropy.} $\Ss(\rho)=h(0.854)\approx0.854\times0.228+0.146\times2.776\approx0.195+0.405=0.601$ bits.
    \item \emph{Comparison.} The ensemble probabilities are $(\tfrac12,\tfrac12)$, with Shannon entropy $1$ bit. So $\Ss(\rho)=0.601<1$: because $\ket0$ and $\ket+$ are not orthogonal, they overlap, and the mixture is \emph{less} random than a fair coin choosing between two perfectly distinguishable states.
  \end{steps}
\end{solution}

\subsection{Properties}

\begin{proposition}[Basic properties]\label{prop:vn-props}
  Let $\rho,\sigma$ be density operators.
  \begin{enumerate}
    \item[(a)] $0\le\Ss(\rho)\le\log d$; $\Ss(\rho)=0$ iff $\rho$ is pure; $\Ss(\rho)=\log d$ iff $\rho=\Id/d$.
    \item[(b)] \emph{Unitary invariance:} $\Ss(U\rho U^\dagger)=\Ss(\rho)$.
    \item[(c)] \emph{Additivity:} $\Ss(\rho\otimes\sigma)=\Ss(\rho)+\Ss(\sigma)$.
    \item[(d)] \emph{Orthogonal mixtures:} if $\rho_i$ have mutually orthogonal supports, $\Ss\bigl(\sum_ip_i\rho_i\bigr)=\Hs(p)+\sum_ip_i\Ss(\rho_i)$.
  \end{enumerate}
\end{proposition}
\begin{proof}
  \begin{steps}
    \item (a) $\Ss(\rho)$ is the Shannon entropy of the spectrum, so Proposition~\ref{prop:shannon-props} applies: it is $0$ iff one eigenvalue is $1$, i.e.\ $\rho$ is a pure state (Section~\ref{sec:density}), and equals $\log d$ iff all eigenvalues equal $\tfrac1d$, i.e.\ $\rho=\Id/d$.
    \item (b) $U\rho U^\dagger$ has the same eigenvalues as $\rho$ (its eigenvectors are $U\ket{v_i}$).
    \item (c) The eigenvalues of $\rho\otimes\sigma$ are the products $\lambda_i\mu_j$, so $\Ss=-\sum_{i,j}\lambda_i\mu_j(\log\lambda_i+\log\mu_j)=-\sum_i\lambda_i\log\lambda_i\sum_j\mu_j-\sum_j\mu_j\log\mu_j\sum_i\lambda_i$, which is $\Ss(\rho)+\Ss(\sigma)$ because $\sum\lambda_i=\sum\mu_j=1$.
    \item (d) Diagonalise each $\rho_i=\sum_j\lambda^{(i)}_j\ket{v^{(i)}_j}\!\bra{v^{(i)}_j}$. Orthogonal supports mean all the vectors $\ket{v^{(i)}_j}$ together are orthonormal, so they are the eigenvectors of $\sum_ip_i\rho_i$, with eigenvalues $p_i\lambda^{(i)}_j$. Then
      \[
        \Ss=-\sum_{i,j}p_i\lambda^{(i)}_j\bigl(\log p_i+\log\lambda^{(i)}_j\bigr)=-\sum_ip_i\log p_i+\sum_ip_i\Bigl(-\sum_j\lambda^{(i)}_j\log\lambda^{(i)}_j\Bigr). \qedhere
      \]
  \end{steps}
\end{proof}

Two further properties, which we state without proof here and which follow from relative entropy (Section~\ref{sec:relent}), are \emph{concavity}, $\Ss\bigl(\sum_ip_i\rho_i\bigr)\ge\sum_ip_i\Ss(\rho_i)$ (mixing increases entropy), and the general bound $\Ss\bigl(\sum_ip_i\rho_i\bigr)\le\Hs(p)+\sum_ip_i\Ss(\rho_i)$, with equality exactly in the orthogonal case (d).

\subsection{Entropy of entanglement}

Now the most striking application. A pure state of two systems has no uncertainty at all, $\Ss(\rho^{AB})=0$, but its parts can be very uncertain.

\begin{theorem}[Entropy of entanglement]\label{thm:entanglement-entropy}
  If $\ket\psi$ is a pure state of $AB$, then $\Ss(\rho^A)=\Ss(\rho^B)$. This common value, which we write $E(\ket\psi)$, is the \term{entropy of entanglement} of $\ket\psi$. If $\ket\psi=\sum_i\lambda_i\ket{i_A}\ket{i_B}$ is a Schmidt decomposition, $E(\ket\psi)=\Hs(\lambda_1^2,\dots,\lambda_d^2)$.
\end{theorem}
\begin{proof}
  \begin{steps}
    \item By Corollary~\ref{cor:schmidt}(a), $\rho^A$ and $\rho^B$ have the same non-zero eigenvalues $\lambda_i^2$.
    \item The entropy is determined by the non-zero eigenvalues (since $0\log0=0$), so $\Ss(\rho^A)=\Ss(\rho^B)=-\sum_i\lambda_i^2\log\lambda_i^2$. \qedhere
  \end{steps}
\end{proof}

By Proposition~\ref{prop:vn-props}(a), $E=0$ exactly for product states, and $E=\log d$ exactly for the maximally entangled states. A state with $E=1$ in two qubits is said to carry one \term{ebit} of entanglement.

\begin{example}\label{ex:ent-entropy}
  Compute the entanglement entropy of $\ket{\Phi^+}$, of $\ket+\ket0$, and of $\frac1{\sqrt3}(\ket{00}+\ket{01}+\ket{10})$.
\end{example}
\begin{solution}
  \begin{steps}
    \item $\ket{\Phi^+}$: the reduced state is $\Id/2$ (Example~\ref{ex:reduced-bell}), so $E=\Ss(\Id/2)=1$ ebit: maximal.
    \item $\ket+\ket0$: a product state, $\rho^A=\ket+\!\bra+$ is pure, so $E=0$.
    \item The third state: the eigenvalues of $\rho^A$ are $\tfrac{3\pm\sqrt5}{6}\approx0.873,\ 0.127$ (Example~\ref{ex:schmidt}), so $E=h(0.127)\approx0.127\times2.975+0.873\times0.196\approx0.378+0.171=0.550$ ebits.
  \end{steps}
  The third state is entangled, but a little over half as much as a Bell state.
\end{solution}

\subsection{Quantum entropy is not classical entropy}

Compare a Bell state with its classical counterpart, the mixture of $\ket{00}$ and $\ket{11}$.

\begin{example}[Bell state against a classical mixture]\label{ex:bell-vs-mixture}
  Let $\rho_1=\ket{\Phi^+}\!\bra{\Phi^+}$ and $\rho_2=\tfrac12\ket{00}\!\bra{00}+\tfrac12\ket{11}\!\bra{11}$. Compute $\Ss(AB)$, $\Ss(A)$, $\Ss(B)$ and the mutual information $I(A\!:\!B)=\Ss(A)+\Ss(B)-\Ss(AB)$ for both.
\end{example}
\begin{solution}
  \begin{steps}
    \item \emph{The Bell state.} $\rho_1$ is pure, so $\Ss(AB)=0$. Each qubit alone is maximally mixed, $\Ss(A)=\Ss(B)=1$. Hence $I=1+1-0=2$ bits.
    \item \emph{The mixture.} $\rho_2$ has eigenvalues $\tfrac12,\tfrac12,0,0$, so $\Ss(AB)=1$. Its reduced states are $\tfrac12\ket0\!\bra0+\tfrac12\ket1\!\bra1=\Id/2$ again, so $\Ss(A)=\Ss(B)=1$. Hence $I=1+1-1=1$ bit.
  \end{steps}
  Both states give the same marginals and the same $Z$ correlations, but the Bell state has \emph{twice} the mutual information. And $\Ss(AB)=0<\Ss(A)=1$: the whole is less uncertain than its part. For classical variables this is impossible, since $\Hs(X,Y)\ge\Hs(X)$ always. It is another face of entanglement.
\end{solution}

The general inequality that remains true is \term{subadditivity}, $\Ss(AB)\le\Ss(A)+\Ss(B)$, which we prove in Section~\ref{sec:relent}; the Araki--Lieb inequality $\abs{\Ss(A)-\Ss(B)}\le\Ss(AB)$ is its companion.

\begin{lookahead}
  Two operational statements make this concrete. \emph{Compression:} $n$ copies of a state $\rho$ can be stored in about $n\,\Ss(\rho)$ qubits (Schumacher), the exact analogue of Theorem~\ref{thm:source-coding}. \emph{Communication:} by Holevo's theorem, a qubit can carry at most one bit of classical information \emph{if used alone}, whatever ensemble is used, since
  the accessible information is at most $\Ss(\rho)-\sum_ip_i\Ss(\rho_i)\le\log d$. Superdense coding (Section~\ref{sec:superdense}) does not violate this: it uses the entangled partner Bob already holds, and the Bell state of Example~\ref{ex:bell-vs-mixture} is precisely the $I(A\!:\!B)=2$ that makes two bits per qubit possible.
\end{lookahead}

\begin{exerciseset}
  \exercise{1} Compute the von Neumann entropy of $\ket+\!\bra+$, $\tfrac12\Id$, $\tfrac13\ket0\!\bra0+\tfrac23\ket1\!\bra1$ and $\diag(\tfrac12,\tfrac14,\tfrac14)$.

  \exercise{1} A qubit has Bloch vector $\vec r=(0,0,\tfrac35)$. Compute its eigenvalues and its entropy. Repeat for $\abs{\vec r}=0.8$ and $\abs{\vec r}=0.9$ (use Figure~\ref{fig:entropy-ball} as a check).

  \exercise{2} Prove parts (a)--(c) of Proposition~\ref{prop:vn-props} for a $2$-level system, and verify (d) for $\rho_1=\ket0\!\bra0$, $\rho_2=\ket1\!\bra1$, $p=(\tfrac13,\tfrac23)$.

  \exercise{2} Compute the entropy of $\tfrac12\ket0\!\bra0+\tfrac12\ket{+}\!\bra{+}$ and of $\tfrac12\ket0\!\bra0+\tfrac12\ket{1}\!\bra{1}$, and explain why they differ although both mix two states with probabilities $(\tfrac12,\tfrac12)$.

  \exercise{2} Show that rotating a qubit state on the Bloch sphere does not change its entropy. (Hint: Proposition~\ref{prop:vn-props}(b).)

  \exercise{2} Compute the entanglement entropy of $\bell xy$, of $\tfrac1{\sqrt5}(\ket{00}+2\ket{11})$ and of $\frac1{\sqrt2}(\ket{00}+\ket{01})$.

  \exercise{2} For the state $\frac12(\ket{00}+\ket{01}+\ket{10}-\ket{11})$ find the Schmidt coefficients and the entanglement entropy.

  \exercise{3} Show that every pure state of two qubits has $0\le E\le1$, with $E=1$ iff it is maximally entangled, and express $E$ in terms of $\abs{\det C}$ (the concurrence): for a two-qubit state with Schmidt coefficients $\lambda_1,\lambda_2$ one has $2\abs{\det C}=2\lambda_1\lambda_2$.

  \exercise{2} Compute $\Ss(\rho^A)$, $\Ss(\rho^B)$, $\Ss(\rho^{AB})$ and $I(A\!:\!B)$ for $\ket0\ket+$ and for $\bell11$.

  \exercise{3} Prove $\Ss\bigl(\tfrac12(\rho_1+\rho_2)\bigr)\ge\tfrac12\Ss(\rho_1)+\tfrac12\Ss(\rho_2)$ (concavity) for commuting $\rho_1,\rho_2$, using the concavity of the Shannon entropy.

  \exercise{3} A source emits $\ket0$ with probability $\tfrac12$ and $\ket+$ with probability $\tfrac12$. Using Example~\ref{ex:vn-mixture}, how many qubits per symbol are needed to store a long message, asymptotically, according to Schumacher's theorem? Compare with the $1$ bit one would need if the two states were orthogonal.
\end{exerciseset}

%======================================================================
\section{Relative entropy}\label{sec:relent}
%======================================================================

\begin{objectives}
  You will compute the classical and quantum relative entropies, prove Gibbs' and Klein's inequalities, and use them to prove the main properties of entropies: the non-negativity of the mutual information, subadditivity, concavity, and the fact that measuring in a basis increases entropy.
\end{objectives}

\begin{recall}
  The von Neumann entropy (Section~\ref{sec:vonneumann}) and Shannon entropy, Jensen's inequality (Section~\ref{sec:shannon}); the partial trace and $\Tr((M\otimes\Id)\rho^{AB})=\Tr(M\rho^A)$ (Section~\ref{sec:density}); projectors and the functional calculus (Section~\ref{sec:operators}).
\end{recall}

How different are two probability distributions? Two quantum states? The entropy of a single distribution cannot answer this, but there is a natural asymmetric measure of how surprised we are to see data generated by $p$ when we believed it was generated by $q$: the relative entropy. It is also the tool that gives the cleanest proofs of nearly every inequality in this chapter.

\subsection{The classical relative entropy}

\begin{definition}[Relative entropy of distributions]\label{def:kl}
  For probability distributions $p,q$ on the same finite set, the \term{relative entropy} (or Kullback--Leibler divergence) is
  \[
    \KL pq=\sum_xp(x)\log\frac{p(x)}{q(x)},
  \]
  with the conventions $0\log\frac0q=0$ and $p\log\frac p0=+\infty$ for $p>0$.
\end{definition}

It is the average, under $p$, of the \emph{extra} surprise $\log\frac{p}{q}$ of $q$ over $p$. It is not symmetric, and it is not a distance, but it is a faithful measure of discrepancy, by the following inequality.

\begin{theorem}[Gibbs' inequality]\label{thm:gibbs}
  For all distributions $p,q$, $\ \KL pq\ge0$, and $\KL pq=0$ iff $p=q$. More generally the inequality holds if $q$ is only a \emph{sub-normalised} distribution ($q(x)\ge0$, $\sum_xq(x)\le1$).
\end{theorem}
\begin{proof}
  If $p(x)>0$ for some $x$ with $q(x)=0$, then $\KL pq=\infty$ and there is nothing to prove, so assume $\supp p\subseteq\supp q$.
  \begin{steps}
    \item Use the inequality $\ln t\le t-1$, with equality iff $t=1$.
    \item \emph{Apply it to $t=q(x)/p(x)$ for each $x$ with $p(x)>0$:} $-\KL pq=\frac1{\ln2}\sum_{p(x)>0}p(x)\ln\frac{q(x)}{p(x)}\le\frac1{\ln2}\sum_{p(x)>0}p(x)\Bigl(\frac{q(x)}{p(x)}-1\Bigr)$.
    \item \emph{Simplify:} the right-hand side is $\frac1{\ln2}\bigl(\sum_{p(x)>0}q(x)-1\bigr)\le0$, since $\sum q\le1$.
    \item \emph{Equality:} we need $q(x)=p(x)$ whenever $p(x)>0$ and $\sum_{p(x)>0}q(x)=1$, which together force $q=p$. \qedhere
  \end{steps}
\end{proof}

Gibbs' inequality gives two earlier results at once, and a third one that we had postponed.

\begin{corollary}\label{cor:gibbs-consequences}
  \begin{enumerate}
    \item[(a)] $\Hs(p)\le\log n$, with equality iff $p$ is uniform.
    \item[(b)] $I(X\!:\!Y)=\KL{p_{XY}}{p_Xp_Y}\ge0$, with equality iff $X,Y$ are independent.
    \item[(c)] Conditioning reduces entropy: $\Hs(Y\mid X)\le\Hs(Y)$.
  \end{enumerate}
\end{corollary}
\begin{proof}
  \begin{steps}
    \item (a) Take $q$ uniform: $\KL pq=\sum_xp(x)\log(np(x))=\log n-\Hs(p)\ge0$.
    \item (b) Expanding, $\KL{p_{XY}}{p_Xp_Y}=\sum p(x,y)\bigl[\log p(x,y)-\log p_X(x)-\log p_Y(y)\bigr]=-\Hs(X,Y)+\Hs(X)+\Hs(Y)=I(X\!:\!Y)$. Gibbs gives $\ge0$, with equality iff $p(x,y)=p_X(x)p_Y(y)$.
    \item (c) $\Hs(Y)-\Hs(Y\mid X)=I(X\!:\!Y)\ge0$. \qedhere
  \end{steps}
\end{proof}

\begin{example}\label{ex:kl}
  Compute $\KL pq$ and $\KL qp$ for $p=(\tfrac12,\tfrac12)$ and $q=(\tfrac14,\tfrac34)$.
\end{example}
\begin{solution}
  \begin{steps}
    \item $\KL pq=\tfrac12\log\frac{1/2}{1/4}+\tfrac12\log\frac{1/2}{3/4}=\tfrac12(1)+\tfrac12(-0.585)\approx0.208$ bits.
    \item $\KL qp=\tfrac14\log\frac{1/4}{1/2}+\tfrac34\log\frac{3/4}{1/2}=\tfrac14(-1)+\tfrac34(0.585)\approx0.189$ bits.
  \end{steps}
  Both are positive, as Gibbs says, but they differ: relative entropy is not symmetric.
\end{solution}

\begin{lookahead}
  Relative entropy is the exact measure of how well a test can tell $p$ from $q$. By Stein's lemma, given $n$ independent samples, any test that wrongly accepts $q$ with probability bounded away from $1$ when the truth is $p$ can make its other error shrink at best like $2^{-n\,\KL pq}$. The quantum version, below, plays the same role for distinguishing quantum states, and extends the work of Section~\ref{sec:measurement} on distinguishing states with a single measurement.
\end{lookahead}

\subsection{The quantum relative entropy}

\begin{definition}[Quantum relative entropy]\label{def:qre}
  For density operators $\rho,\sigma$, the \term{relative entropy} of $\rho$ with respect to $\sigma$ is
  \[
    \Srel\rho\sigma=\Tr\bigl(\rho\log\rho\bigr)-\Tr\bigl(\rho\log\sigma\bigr)\quad\text{if }\supp\rho\subseteq\supp\sigma,\qquad\text{and }+\infty\text{ otherwise}.
  \]
  Here $\log\sigma$ is the operator logarithm on the support of $\sigma$ (Definition~\ref{def:functional}).
\end{definition}

If $\rho$ and $\sigma$ commute, they share an eigenbasis and $\Srel\rho\sigma$ is just the classical $\KL pq$ for the two lists of eigenvalues. The quantum case adds the possibility that they do \emph{not} commute.

\begin{example}\label{ex:qre}
  Compute $\Srel\rho\sigma$ for: (a) $\rho=\ket0\!\bra0$, $\sigma=\Id/2$; (b) $\rho=\ket0\!\bra0$, $\sigma=\ket1\!\bra1$; (c) $\rho=\diag(\tfrac34,\tfrac14)$, $\sigma=\Id/2$; (d) $\rho=\diag(\tfrac34,\tfrac14)$ and $\sigma=\tfrac34\ket+\!\bra++\tfrac14\ket-\!\bra-$.
\end{example}
\begin{solution}
  \begin{steps}
    \item (a) $\Tr\rho\log\rho=0$ and $\log\sigma=-\Id$ (in bits), so $\Srel\rho\sigma=0-(-1)=1$ bit.
    \item (b) $\supp\rho=\Span\{\ket0\}\not\subseteq\supp\sigma=\Span\{\ket1\}$, so $\Srel\rho\sigma=+\infty$: a state that is certainly $\ket0$ is infinitely surprising if we believed $\ket1$.
    \item (c) $\Srel\rho\sigma=-\Ss(\rho)-\Tr\rho\log\sigma=-0.811+1=0.189$ bits, which is $\log2-\Ss(\rho)$.
    \item (d) The operators do not commute. We use $\log\sigma=\log\tfrac34\,\ket+\!\bra++\log\tfrac14\,\ket-\!\bra-$, so $\Tr\rho\log\sigma=\log\tfrac34\bra+\rho\ket++\log\tfrac14\bra-\rho\ket-$. Since $\bra\pm\rho\ket\pm=\tfrac12(\tfrac34+\tfrac14)=\tfrac12$, this is $\tfrac12\log\tfrac34+\tfrac12\log\tfrac14\approx-0.208-1=-1.208$. Hence $\Srel\rho\sigma=-0.811+1.208=0.396$ bits.
  \end{steps}
\end{solution}

The fundamental result is the quantum Gibbs inequality, due to Klein.

\begin{theorem}[Klein's inequality]\label{thm:klein}
  $\Srel\rho\sigma\ge0$, with equality iff $\rho=\sigma$.
\end{theorem}
\begin{proof}
  Assume $\supp\rho\subseteq\supp\sigma$ (otherwise the left side is $+\infty$). Write $\rho=\sum_ip_i\ket i\!\bra i$ and $\sigma=\sum_jq_j\ket{j'}\!\bra{j'}$ with two orthonormal bases, in general different. Let $P_{ij}=\abs{\braket i{j'}}^2$.
  \begin{steps}
    \item \emph{The matrix $P$.} $\sum_jP_{ij}=1$ for each $i$ (completeness of $\{\ket{j'}\}$), and $\sum_iP_{ij}=1$ for each $j$. So the rows and columns of $P$ are probability distributions: $P$ is \emph{doubly stochastic}.
    \item \emph{Rewrite the relative entropy.} $\Tr\rho\log\sigma=\sum_ip_i\bra i\log\sigma\ket i=\sum_ip_i\sum_jP_{ij}\log q_j$ (sums over $q_j>0$; the terms with $p_i>0$ only involve this part of $\sigma$ because $\supp\rho\subseteq\supp\sigma$). Hence
      \[
        \Srel\rho\sigma=\sum_ip_i\log p_i-\sum_ip_i\sum_jP_{ij}\log q_j .
      \]
    \item \emph{Jensen.} For fixed $i$, $(P_{ij})_j$ is a distribution and $\log$ is concave, so $\sum_jP_{ij}\log q_j\le\log\sum_jP_{ij}q_j=\log r_i$, where $r_i=\sum_jP_{ij}q_j$.
    \item \emph{A classical divergence appears.} Therefore $\Srel\rho\sigma\ge\sum_ip_i\log\frac{p_i}{r_i}=\KL pr$.
    \item \emph{$r$ is sub-normalised.} $\sum_ir_i=\sum_jq_j\sum_iP_{ij}=\sum_jq_j=1$ (the restriction to $q_j>0$ only makes it smaller, never larger than $1$).
    \item \emph{Conclude.} By Gibbs' inequality (Theorem~\ref{thm:gibbs}), $\KL pr\ge0$, so $\Srel\rho\sigma\ge0$.
    \item \emph{Equality.} If $\Srel\rho\sigma=0$, then equality holds in Jensen: for each $i$ with $p_i>0$, $q_j$ is constant on those $j$ with $P_{ij}>0$, equal to $r_i$; and $p=r$. Hence $\sigma\ket i=\sum_jq_j\braket{j'}{i}\ket{j'}=p_i\ket i$ for each $i$ with $p_i>0$. These eigenvectors already account for trace $\sum_ip_i=1$ of $\sigma$, so $\sigma$ has no further non-zero eigenvalue and $\sigma=\rho$. \qedhere
  \end{steps}
\end{proof}

\subsection{Consequences}

Klein's inequality quickly produces the properties promised in Section~\ref{sec:vonneumann}.

\begin{corollary}[Properties of the von Neumann entropy]\label{cor:vn-consequences}
  \begin{enumerate}
    \item[(a)] $\Ss(\rho)\le\log d$, with equality iff $\rho=\Id/d$.
    \item[(b)] \emph{Subadditivity:} $\Ss(\rho^{AB})\le\Ss(\rho^A)+\Ss(\rho^B)$, with equality iff $\rho^{AB}=\rho^A\otimes\rho^B$. Equivalently, the quantum mutual information $I(A\!:\!B)=\Ss(\rho^A)+\Ss(\rho^B)-\Ss(\rho^{AB})=\Srel{\rho^{AB}}{\rho^A\otimes\rho^B}\ge0$.
    \item[(c)] \emph{Concavity:} $\Ss\bigl(\sum_ip_i\rho_i\bigr)\ge\sum_ip_i\Ss(\rho_i)$.
    \item[(d)] \emph{Measuring in a basis increases entropy:} if $\mathcal D(\rho)=\sum_i\ket i\!\bra i\rho\ket i\!\bra i$ is the state after a non-selective measurement in an orthonormal basis $\{\ket i\}$, then $\Ss(\mathcal D(\rho))=\Ss(\rho)+\Srel\rho{\mathcal D(\rho)}\ge\Ss(\rho)$.
  \end{enumerate}
\end{corollary}
\begin{proof}
  \begin{steps}
    \item (a) Take $\sigma=\Id/d$: $\Srel\rho{\Id/d}=\Tr\rho\log\rho-\Tr\rho\log\frac{\Id}d=-\Ss(\rho)+\log d\ge0$, with equality iff $\rho=\Id/d$.
    \item (b) \emph{The logarithm of a product.} If $\rho^A=\sum_ap_a\ket a\!\bra a$ and $\rho^B=\sum_bq_b\ket b\!\bra b$ then $\rho^A\otimes\rho^B$ has eigenvectors $\ket a\ket b$ with eigenvalues $p_aq_b$, so $\log(\rho^A\otimes\rho^B)=\log\rho^A\otimes\Id+\Id\otimes\log\rho^B$. \emph{Take the expectation} in $\rho^{AB}$ using Proposition~\ref{prop:reduced}: $\Tr\bigl(\rho^{AB}(\log\rho^A\otimes\Id)\bigr)=\Tr\rho^A\log\rho^A=-\Ss(\rho^A)$, and likewise for $B$. Hence $\Srel{\rho^{AB}}{\rho^A\otimes\rho^B}=-\Ss(\rho^{AB})+\Ss(\rho^A)+\Ss(\rho^B)$, which is $\ge0$ by Klein, with equality iff $\rho^{AB}=\rho^A\otimes\rho^B$.
    \item (c) Let $\bar\rho=\sum_ip_i\rho_i$. Since $\supp\rho_i\subseteq\supp\bar\rho$, $\sum_ip_i\Srel{\rho_i}{\bar\rho}=\sum_ip_i\bigl(-\Ss(\rho_i)\bigr)-\Tr\bar\rho\log\bar\rho=\Ss(\bar\rho)-\sum_ip_i\Ss(\rho_i)$, and each term of the left side is $\ge0$.
    \item (d) $\log\mathcal D(\rho)$ is diagonal in $\{\ket i\}$, and for any diagonal $D$ one has $\Tr(\rho D)=\Tr(\mathcal D(\rho)D)$. Hence $\Tr\rho\log\mathcal D(\rho)=\Tr\mathcal D(\rho)\log\mathcal D(\rho)=-\Ss(\mathcal D(\rho))$. So $\Srel\rho{\mathcal D(\rho)}=-\Ss(\rho)+\Ss(\mathcal D(\rho))\ge0$. \qedhere
  \end{steps}
\end{proof}

\begin{example}\label{ex:entropy-increase}
  Verify (d) for $\rho=\ket+\!\bra+$ with the computational basis.
\end{example}
\begin{solution}
  \begin{steps}
    \item \emph{Before:} $\rho$ is pure, so $\Ss(\rho)=0$.
    \item \emph{After:} $\mathcal D(\rho)=\tfrac12\ket0\!\bra0+\tfrac12\ket1\!\bra1=\Id/2$, with $\Ss=1$.
    \item \emph{The relative entropy:} $\Srel\rho{\mathcal D(\rho)}=1-0=1$, consistent with the identity of (d).
  \end{steps}
  A projective measurement that we forget the result of destroys the coherence between $\ket0$ and $\ket1$ and raises the entropy from $0$ to its maximum.
\end{solution}

\begin{example}[Distinguishability can only be lost by discarding]\label{ex:monotone}
  Compare $\Srel{\rho^{AB}}{\sigma^{AB}}$ for $\rho^{AB}=\ket{\Phi^+}\!\bra{\Phi^+}$, $\sigma^{AB}=\Id/4$ with the relative entropy of the reduced states.
\end{example}
\begin{solution}
  \begin{steps}
    \item \emph{Joint:} $\Srel{\rho^{AB}}{\Id/4}=\log4-\Ss(\rho^{AB})=2-0=2$ bits.
    \item \emph{Reduced:} $\rho^A=\Id/2=\sigma^A$, so $\Srel{\rho^A}{\sigma^A}=0$.
  \end{steps}
  Two bits of distinguishability between the global states vanish when we look at one qubit only: all of it lived in the correlations. In general $\Srel{\rho^A}{\sigma^A}\le\Srel{\rho^{AB}}{\sigma^{AB}}$ (\emph{monotonicity under partial trace}), a deep theorem of Lieb, Ruskai and Lindblad, which gives the strong subadditivity of the entropy and says that no physical operation can increase distinguishability. It is proved in~\cite{wilde2017}.
\end{solution}

\begin{lookahead}
  Relative entropy is the unifying quantity. The mutual information is a relative entropy to a product; the entanglement of a state can be quantified by its relative entropy to the closest unentangled state; and in thermodynamics, with $\sigma=e^{-\beta H}/Z$ the Gibbs state, $\Srel\rho\sigma=\beta\bigl(F(\rho)-F(\sigma)\bigr)$ is the excess free energy of $\rho$ (Section~\ref{sec:blackbody} uses this identity to show that the Gibbs state has the largest entropy at fixed energy). It is also the exponent that governs the optimal discrimination of two quantum states from many copies (the quantum Stein lemma), and the natural next step after the one-shot discrimination of Section~\ref{sec:measurement}. See Vedral's review~\cite{vedral2002}.
\end{lookahead}

\begin{exerciseset}
  \exercise{1} Compute $\KL pq$ and $\KL qp$ for $p=(0.9,0.1)$, $q=(0.5,0.5)$, and for $p=(\tfrac12,\tfrac12,0)$, $q=(\tfrac13,\tfrac13,\tfrac13)$ and the reverse.

  \exercise{2} Prove $\ln t\le t-1$ for $t>0$ (study $g(t)=t-1-\ln t$), with equality only at $t=1$, and use it to give a one-line proof that $\KL pq\ge0$ for sub-normalised $q$.

  \exercise{2} Show that the mutual information can be written $I(X\!:\!Y)=\sum_{x,y}p(x,y)\log\frac{p(x,y)}{p_X(x)p_Y(y)}$, and compute it for the joint distribution of Exercise~\ref{ex:joint-dist}.

  \exercise{2} For the binary distributions $p=(a,1-a)$ and $q=(b,1-b)$ write $\KL pq$ explicitly, and show it vanishes only for $a=b$.

  \exercise{1} Compute $\Srel\rho\sigma$ for $\rho=\ket+\!\bra+$ and $\sigma=\Id/2$, for $\rho=\Id/2$ and $\sigma=\ket0\!\bra0$, and for $\rho=\diag(\tfrac12,\tfrac12)$, $\sigma=\diag(\tfrac14,\tfrac34)$.

  \exercise{2} Repeat Example~\ref{ex:qre}(d) with $\sigma=\tfrac12\ket+\!\bra++\tfrac12\ket-\!\bra-=\Id/2$ and with $\sigma=\tfrac{9}{10}\ket+\!\bra++\tfrac1{10}\ket-\!\bra-$, and compare.

  \exercise{2} Show that $\Srel\rho\sigma$ is not symmetric by computing it both ways for $\rho=\diag(\tfrac34,\tfrac14)$ and $\sigma=\diag(\tfrac12,\tfrac12)$.

  \exercise{2} Verify Corollary~\ref{cor:vn-consequences}(c) (concavity) for $\rho_1=\ket0\!\bra0$, $\rho_2=\ket+\!\bra+$ and $p=(\tfrac12,\tfrac12)$ by computing both sides explicitly.

  \exercise{2} Use the formula $\log(\rho\otimes\sigma)=\log\rho\otimes\Id+\Id\otimes\log\sigma$ to prove that $\Srel{\rho_1\otimes\rho_2}{\sigma_1\otimes\sigma_2}=\Srel{\rho_1}{\sigma_1}+\Srel{\rho_2}{\sigma_2}$.

  \exercise{3} Show that $\Srel{U\rho U^\dagger}{U\sigma U^\dagger}=\Srel\rho\sigma$ for every unitary $U$ (unitary invariance), and deduce that distinguishability is the same in any basis.

  \exercise{2} Verify the subadditivity inequality for $\bell00$, for $\ket+\ket0$, and for the state $\tfrac1{\sqrt3}(\ket{00}+\ket{01}+\ket{10})$ (use Example~\ref{ex:ent-entropy}).

  \exercise{3} Show that $I(A\!:\!B)\le2\min\{\Ss(\rho^A),\Ss(\rho^B)\}$ for a pure state of $AB$, and that equality holds for all pure states. For a general mixed state the bound is $I\le2\min\{\log d_A,\log d_B\}$.

  \exercise{3} For $\ket\psi=\cos\frac\theta2\ket0+\sin\frac\theta2\ket1$, compute $\Ss(\mathcal D(\ket\psi\!\bra\psi))$ for the dephasing in the computational basis, as a function of $\theta$, and check that it is $\ge0=\Ss(\ket\psi\!\bra\psi)$.

  \exercise{3} Let $\rho_0=\ket0\!\bra0$ and $\rho_1=\ket+\!\bra+$ be sent with probability $\tfrac12$ each. Compute the \emph{Holevo quantity} $\chi=\Ss(\bar\rho)-\tfrac12\Ss(\rho_0)-\tfrac12\Ss(\rho_1)$ and compare it with $1$ bit. (It bounds the information that Bob can extract about the bit sent.)
\end{exerciseset}

\begin{furtherreading}
  \item Billingsley~\cite{billingsley1995}: the standard rigorous introduction to probability measures; for the measure-theoretic background that we used lightly in Section~\ref{sec:prob}.
  \item Stein and Shakarchi~\cite{steinshakarchi2003}: a clear and complete first course on Fourier series and transforms, with the proofs we omitted (completeness, inversion, Plancherel).
  \item Shannon's original paper~\cite{shannon1948} and Cover and Thomas~\cite{cover2006}: the foundation and the standard textbook of classical information theory.
  \item Wilde~\cite{wilde2017} and Nielsen and Chuang~\cite{nielsenchuang2010} (Chapters~11 and~12): quantum entropies, strong subadditivity, Holevo's theorem and Schumacher compression.
  \item Vedral~\cite{vedral2002}: a review of the many roles of the relative entropy, from thermodynamics to entanglement.
\end{furtherreading}

